%====================================================
% Branch X — Information Geometry of Semantic Fields
%====================================================
\documentclass[11pt]{article}
\usepackage{amsmath}
% ---------- Encoding & Fonts ----------
\usepackage[T1]{fontenc}
\usepackage[utf8]{inputenc}
\usepackage{microtype}
\usepackage{lmodern}
\usepackage{newunicodechar}
\newunicodechar{→}{\textrightarrow{}}
\newunicodechar{∼}{\textasciitilde{}}
\newunicodechar{≤}{\ensuremath{\leq}}
\newunicodechar{−}{\ensuremath{-}}
\newunicodechar{≈}{\ensuremath{\approx}}

% ---------- Page Layout ----------
\usepackage[margin=1in]{geometry}
\usepackage{setspace}
\setstretch{1.08}

% ---------- Math & Symbols ----------
\usepackage{amsmath,amssymb,amsfonts,amsthm,bm}
\usepackage{mathtools}
\usepackage{enumitem}
\usepackage{tabularx}
\usepackage{siunitx}
\sisetup{detect-all=true}
\DeclareMathOperator*{\argmin}{arg\,min}
\DeclareMathOperator*{\argmax}{arg\,max}
\DeclareMathOperator{\tr}{tr}
\DeclareMathOperator{\diag}{diag}
\DeclareMathOperator{\KL}{KL}
\newcommand{\E}{\mathbb{E}}
\newcommand{\Prob}{\mathbb{P}}
\newcommand{\Var}{\mathrm{Var}}
\newcommand{\Cov}{\mathrm{Cov}}
\newcommand{\norm}[1]{\left\lVert #1 \right\rVert}
\newcommand{\abs}[1]{\left| #1 \right|}
\newcommand{\RR}{\mathbb{R}}
\newcommand{\NN}{\mathbb{N}}
\newcommand{\gF}{g^{\mathrm{F}}}
\newcommand{\DKL}{D_{\mathrm{KL}}}
\newcommand{\SamTrans}{\mathsf{SamTrans}}
\newcommand{\SamGeom}{\mathsf{SamGeom}}
\newcommand{\SamInfo}{\mathsf{SamInfo}}
\newcommand{\MSR}{\mathsf{MSR}}

% ---------- Graphics & Tables ----------
\usepackage{graphicx}
\usepackage{booktabs}
\usepackage{multirow}
\usepackage{adjustbox}

% ---------- Framed Boxes ----------
\usepackage{mdframed}
\usepackage{xcolor}

% ---------- Algorithms ----------
\usepackage{algorithm}
\usepackage{algpseudocode}
\usepackage{float}

\makeatletter
\renewcommand{\fps@algorithm}{htbp}
\setcounter{totalnumber}{10}
\setcounter{topnumber}{10}
\setcounter{bottomnumber}{10}
\renewcommand{\topfraction}{0.9}
\renewcommand{\bottomfraction}{0.9}
\renewcommand{\textfraction}{0.1}
\renewcommand{\floatpagefraction}{0.8}
\makeatother
\setlength{\intextsep}{10pt plus 2pt minus 2pt}

% ---------- Links & Clever Refs ----------
\usepackage[hidelinks]{hyperref}
\usepackage[capitalise,nameinlink]{cleveref}

% ---------- Quotes & Lists ----------
\usepackage{csquotes}

\setlist[itemize]{leftmargin=*, itemsep=2pt, topsep=2pt}
\setlist[enumerate]{leftmargin=*, itemsep=2pt, topsep=2pt}

% ---------- Theorem Envs ----------
\newtheorem{definition}{Definition}
\newtheorem{theorem}{Theorem}
\newtheorem{lemma}{Lemma}
\newtheorem{proposition}{Proposition}
\newtheorem{corollary}{Corollary}
\newtheorem{assumption}{Assumption}
\newtheorem{conjecture}{Conjecture}
\theoremstyle{remark}
\newtheorem{remark}{Remark}

% ---------- Title Block ----------
\title{\textbf{Information Geometry of Semantic Fields:}\\[0.25ex]
\textbf{The Fisher Metric as Canonical Davis Metric}\\[0.5ex]
\large Dual Connections, Spectral Capacity, and Measurable Curvature\\[0.25ex]
\large Branch X of the Davis Framework}
\author{Bee Rosa Davis\\
\small \texttt{bee\_davis@alumni.brown.edu}}
\date{}

\begin{document}
\maketitle
\vspace{-1.5ex}

%====================================================
% ABSTRACT
%====================================================
\begin{abstract}
The Davis framework for geometry-first detection posits a Riemannian manifold
$(M, g)$ on which semantic evolution obeys compositional field equations with
explicit curvature capacity bounds, spectral gap constraints, and
renormalization group structure (Branches~I--IX).  Throughout this development,
the metric~$g$ has remained a free parameter: chosen by the engineer, validated
by distortion audits, never derived from first principles.

This paper resolves the metric-selection problem through an identification, not
an invention.  We observe that when the Davis manifold carries a statistical
structure---that is, when points on $M$ parameterize probability
distributions---\v{C}encov's uniqueness theorem (1982) forces $g$ to be the
Fisher information metric $\gF$, up to a positive scalar.  This is the unique
Riemannian metric invariant under sufficient statistics, and it is exactly the
invariance demanded by the Davis framework's representation-independence
principle.

The identification has three consequences.  First, Amari's
$\alpha$-connection family equips $(M, \gF)$ with a one-parameter deformation
$\nabla^{(\alpha)}$ of the Levi-Civita connection, including dual flat
connections $\nabla^{(e)}$ and $\nabla^{(m)}$ on exponential families.  This
duality has no analog in the parent framework and constitutes the genuine
structural extension of this branch.  Second, the parent's curvature capacity
law is obtained from Bonnet--Myers and the Cram\'er--Rao inequality on the Fisher
manifold, with equality on exponential families explaining why softmax
architectures locally saturate capacity bounds.  Third, the Fisher metric is the
Hessian of the log-partition function, which is the free energy of Branch~VI;
the susceptibility matrix of statistical mechanics, the spectral gap of
Branch~IX, the sheaf compatibility condition of Branch~VII, and the RG
monotonicity of Branch~VIII all acquire information-geometric interpretations.

We verify the identification on the finite combinatorial models (Shidoku and
Latin square CSP graphs) from Branch~IX, computing Fisher-weighted Laplacians
and checking cross-branch consistency.  We close with an empirical measurement
program: the Fisher metric is computable from transformer logit distributions,
making the Davis framework's curvature bounds testable on real neural
architectures for the first time.
\end{abstract}

\tableofcontents
\newpage

%====================================================
% IDENTIFICATION TABLE
%====================================================

\begin{table}[htbp]
\centering
\caption{\textbf{Core identification:} Davis framework objects and their
information-geometric counterparts.  Each row shows a parent-framework
construct (left) and its canonical Fisher-metric realization (right).
Entries marked $\dagger$ are new to Branch~X; all others are classical.}
\label{tab:identification}
\begin{tabularx}{\linewidth}{@{}p{0.38\linewidth}p{0.54\linewidth}@{}}
\toprule
\textbf{Davis framework object} & \textbf{Fisher-metric realization} \\
\midrule
Manifold $(M, g)$ &
Statistical manifold $(\mathcal{S}, \gF)$ with Fisher metric \\
\addlinespace[2pt]
Metric $g$ (free parameter) &
$\gF_{ij} = \E[\partial_i \ell \cdot \partial_j \ell]$ (canonical via \v{C}encov) \\
\addlinespace[2pt]
Levi-Civita connection $\nabla^{(0)}$ &
Self-dual $\alpha = 0$ connection (parent's only connection) \\
\addlinespace[2pt]
--- (no parent analog)$^\dagger$ &
$\alpha$-connection family $\nabla^{(\alpha)}$; dual pair $(\nabla^{(e)}, \nabla^{(m)})$ \\
\addlinespace[2pt]
Sectional curvature $K$ &
Statistical curvature (Efron 1975); $\alpha$-curvature tensor \\
\addlinespace[2pt]
Geodesic paths &
Maximum-likelihood trajectories in distribution space \\
\addlinespace[2pt]
Holonomy $\mathrm{Hol}(\gamma)$ &
Information loss around constraint loops (Ambrose--Singer) \\
\addlinespace[2pt]
Tolerance budget $\tau$ &
KL divergence budget: $\sum_\ell D_{\mathrm{KL}}(p_\ell \| p_{\ell+1}) \le \tau$ \\
\addlinespace[2pt]
Distortion $\varepsilon(L)$ &
Fisher-Rao vs.\ Euclidean distortion along completion paths \\
\addlinespace[2pt]
Spectral gap $\mu_1$ &
Fisher spectral gap $\mu_1^{\mathrm{F}}$ (Fisher Laplacian) \\
\addlinespace[2pt]
Capacity $C \le 1/\sup K$ &
Cram\'er--Rao capacity; saturated on exponential families \\
\addlinespace[2pt]
Parallel transport &
$\alpha$-transport of sufficient statistics$^\dagger$ \\
\addlinespace[2pt]
Free energy $\psi = \log Z$ (Branch~VI) &
Log-partition function; $\gF_{ij} = \partial^2\psi / \partial\theta^i\partial\theta^j$ \\
\addlinespace[2pt]
Completion count $N_k$ (Branch~IX) &
Fisher weight $\gF(L_k) = 1/N_k(L_k)$ on CSP graphs \\
\bottomrule
\end{tabularx}
\end{table}


%====================================================
\section{Introduction}\label{sec:introduction}
%====================================================

The Davis framework is a program for geometry-first detection: given a semantic
sameness structure~$S$, one builds a Riemannian manifold $(M, g)$, reads
decisions off geodesic gaps, and bounds error via a compositional budget
(geometry, linkage, calibration, abstention, independence
slack)~\cite{davis2025manifold, davis2025sameness}.  Across nine branches, the
framework has been developed along several mathematical axes: statistical
mechanics (Branch~VI)~\cite{davis2025statmech}, sheaf cohomology
(Branch~VII)~\cite{davis2025sheaf}, renormalization group
(Branch~VIII)~\cite{davis2025rg}, and spectral geometry
(Branch~IX)~\cite{davis2025spectral}.

A persistent gap has remained: the metric~$g$ is always \emph{given}, never
\emph{derived}.  In the Geometry of Sameness~\cite{davis2025sameness}, the
manifold-based semantic realization (MSR) packages feature spaces $\{V_i\}$,
chart maps $\{\psi_i\}$, and a Riemannian metric $g$ into a single object, but
the definition does not specify which metric to choose.  Engineers select~$g$
through contrastive training objectives, smoothness regularizers, or inherited
hyperbolic/spherical structures, and then audit distortion
post~hoc~\cite{davis2025manifold}.

This branch resolves the gap by observing that the answer already exists in
the mathematical literature, and has for forty years.

\subsection{The observation}

Suppose the Davis manifold carries a \emph{statistical structure}: each point
$\theta \in M$ parameterizes a probability distribution $p_\theta$ over
observations.  This is not a restrictive assumption; it is already true in
every concrete instantiation of the framework.  HERALD's antigenic manifold
parameterizes posterior distributions over viral lineages.  VIDAR's identity
manifold parameterizes face-recognition likelihoods.  AMC's translator graph,
when lifted to a manifold via the functor $F_S$, parameterizes
reconstruction-error distributions across modalities.

On statistical manifolds, there is a classical uniqueness result.

\begin{quote}
\emph{\v{C}encov's theorem} (1982): The Fisher information metric is the unique
(up to positive scalar) Riemannian metric on the space of probability
distributions that is invariant under sufficient statistics.
\end{quote}

The Davis framework's representation-independence principle demands that
detection guarantees be invariant under changes of feature-space coordinates
that preserve sufficient statistics of the underlying sameness structure~$S$.
This is precisely the invariance condition in \v{C}encov's theorem.  Therefore:

\begin{quote}
\emph{Within the statistical-manifold regime (axioms IG-A1--A4),
the Fisher information metric is the unique Davis metric, up to scale.}
\end{quote}

This is the core identification of Branch~X.  It is not a new theorem.  It is
a recognition that two independently motivated invariance requirements---one
from detection theory, one from mathematical statistics---select the same
metric, \emph{conditional on the manifold carrying statistical-model structure}.

\subsection{What is new}

While the Fisher metric itself is classical (Rao 1945, \v{C}encov 1982, Amari
1985), three consequences of its identification as the Davis metric
(within the statistical-manifold regime) are new to this framework:

\paragraph{(1) Dual $\alpha$-connections.}
The Fisher metric on a statistical manifold admits Amari's one-parameter family
of $\alpha$-connections $\nabla^{(\alpha)}$~\cite{amari2000}.  The parent
framework uses only the Levi-Civita connection ($\alpha = 0$).  The exponential
connection ($\alpha = +1$) and mixture connection ($\alpha = -1$) form a
conjugate pair that is flat on exponential families.  This duality has no analog
in the parent framework and constitutes the genuine structural extension of
Branch~X.

\paragraph{(2) Curvature law from Cram\'er--Rao.}
A $d$-dimensional curvature capacity bound is proved via
the Bonnet--Myers theorem, Bishop--Gromov volume comparison, and the
Cram\'er--Rao inequality on the Fisher manifold.  The bound
$C_\delta(\Omega) \le c(d)\,\kappa^{-d/2}\,\delta^{-d}$ is explicit in
both dimension and curvature; the parent framework's law $C \lesssim 1/K$
appears as the $d{=}2$ specialization.
On exponential families (such as softmax outputs), the Efron curvature
correction vanishes and the bound is tight.  This provides a structural
explanation for why softmax architectures are locally geometrically efficient.

\paragraph{(3) Measurable curvature.}
The Fisher metric is computable from output distributions.  Given a
transformer's logit distributions $p_{\ell,t}$ at each layer~$\ell$ and
position~$t$, the Fisher information matrix $\gF_\ell$ can be estimated
(via stochastic approximation or low-rank methods such as K-FAC).  This makes
the Davis framework's curvature bounds \emph{empirically testable} on real
neural architectures---a capability that no previous branch provides.

\subsection{Scope and non-goals}

This paper:
\begin{itemize}
  \item identifies the Fisher metric as the unique Davis metric
  within the statistical-manifold regime, via \v{C}encov's theorem;
  \item introduces the $\alpha$-connection family into the Davis framework
  and derives the dual-connection structure on exponential families;
  \item contains the parent curvature capacity law as the $d{=}2$ specialization, via Bonnet--Myers and
  Cram\'er--Rao, with explicit constants, and proves saturation on exponential
  families;
  \item computes Fisher geometry on the finite CSP models (Shidoku, Latin
  square) from Branch~IX as finite-model verification;
  \item proposes a central conjecture (Fisher spectral gap governs completion
  mixing) and an empirical measurement program for real transformers.
\end{itemize}

This paper does \emph{not}:
\begin{itemize}
  \item construct new functors or $\varepsilon$-equivalence chains (those belong
  in a future synthesis paper);
  \item claim new information-geometric results (all core theorems are
  classical; the contribution is identification and assembly);
  \item report empirical results on neural networks (the measurement program
  is proposed, not executed).
\end{itemize}

\subsection{Roadmap}

The paper develops in three acts.

\textbf{Act~I (Identification).}  Section~\ref{sec:statistical-manifolds}
establishes statistical manifolds and the Fisher metric, states \v{C}encov's
uniqueness theorem, and makes the core identification.
Section~\ref{sec:alpha-connections} introduces the $\alpha$-connection family,
proves duality, and derives dual coordinate systems on exponential families.
Section~\ref{sec:curvature-capacity} proves a $d$-dimensional curvature capacity
bound that contains the parent law as the $d{=}2$ specialization,
from Bonnet--Myers and the Cram\'er--Rao inequality, with explicit constants
and a complete proof, and establishes exponential-family saturation.

\textbf{Act~II (Consequences).}  Section~\ref{sec:spectral-capacity} derives
the Fisher spectral capacity and states the central conjecture.
Section~\ref{sec:directions} collects conjectural connections between dual
connections and neural architecture, clearly labeled as directions rather than
results.

\textbf{Act~III (Computation and synthesis).}
Section~\ref{sec:cross-branch} establishes cross-branch synthesis with
Branches~VI--IX.  Section~\ref{sec:finite-models} computes Fisher geometry on
the Shidoku and Latin square graphs.  Section~\ref{sec:empirical} describes the
empirical measurement program.  Section~\ref{sec:limitations} discusses
limitations and open problems.  Appendix~\ref{sec:appendix} provides detailed
proofs.

\subsection{Running examples}\label{subsec:running-examples}

We ground the theory in three systems from earlier papers, which serve as
running examples throughout.

\paragraph{HERALD (viral antigenic drift).}
HERALD~\cite{davis2025manifold} learns a pullback Riemannian manifold on viral
spike protein sequences where geodesic distance approximates antigenic distance.
In the information-geometric view, each point on the HERALD manifold
parameterizes a posterior distribution over compatible viral lineages given
neutralization assay data.  The Fisher metric on this space is the
susceptibility of lineage posteriors to perturbation in antigenic coordinates.

\paragraph{VIDAR (deepfake detection).}
VIDAR~\cite{davis2025manifold} treats face-recognition embeddings as points on
the hypersphere $\mathbb{S}^{d-1}$.  In the information-geometric view, each
embedding parameterizes a von Mises--Fisher distribution over identity
observations.  The Fisher metric of the von Mises--Fisher family on
$\mathbb{S}^{d-1}$ is proportional to the standard round metric, confirming
that VIDAR's inherited geometry is already Fisher-canonical for its output
distribution family.

\paragraph{AMC (telemetry format switching).}
AMC~\cite{davis2025sameness} uses learned format translators between
heterogeneous telemetry formats.  When lifted to a manifold via the functor
$F_S$, each point parameterizes a reconstruction-error distribution across
modalities.  The Fisher metric on this space controls the sensitivity of
cross-format translation to perturbations in the underlying system state.

In each case, the pattern is the same: the Davis manifold already carries a
statistical structure (explicitly or implicitly), and the Fisher metric
canonically instantiates the abstract metric~$g$.

\subsection{Core assumptions}\label{subsec:assumptions}

For reference, we collect the assumptions under which this branch's results
hold.

\begin{samepage}
\begin{mdframed}[linewidth=0.5pt]
\noindent\textbf{Core Assumptions for Branch~X}
\begin{itemize}[leftmargin=*]
  \item \textbf{IG-A1 (Statistical structure).} The Davis manifold $(M, g)$
  carries a statistical structure: each point $\theta \in M$ parameterizes a
  probability distribution $p_\theta$ over an observation space $\mathcal{X}$,
  and the parameterization is identifiable, smooth, and has common support
  (Definition~\ref{def:stat-manifold}).

  \item \textbf{IG-A2 (Fisher regularity).} The Fisher information matrix
  $\gF(\theta)$ is positive definite with uniformly bounded condition number
  $\kappa(\gF) \le \kappa_{\max}$ on the in-regime set $\Omega^{\mathrm{ideal}}$.
  This ensures the Fisher metric is a well-defined Riemannian metric on the
  operational domain.

  \item \textbf{IG-A3 (Exponential family structure).}  For the strongest
  results (dual flatness, Cram\'er--Rao saturation, Pythagorean decomposition),
  the output distribution family $\{p_\theta\}$ is an exponential family.
  Softmax outputs satisfy this assumption.  Results marked ``exponential family
  only'' require IG-A3; all others hold on general statistical manifolds.

  \item \textbf{IG-A4 (Representation-independence).} The Davis metric $g$ is
  required to be invariant under sufficient statistic reductions of the
  observation model.  This is the invariance condition in \v{C}encov's theorem
  and the representation-independence principle of the parent
  framework~\cite{davis2025manifold}.

  \item \textbf{IG-A5 (Bounded path curvature).}  Along benign completion paths
  $\gamma \in \mathcal{P}(L)$, the sectional curvature of $(\mathcal{S}, \gF)$
  is bounded: $|K^{(0)}(\gamma(t))| \le K_{\max}$ for all $t \in [0,1]$.
  This is the Fisher-metric analog of the parent's curvature control
  assumption~A2.

  \item \textbf{IG-A6 (Finite Fisher information).}  The Fisher information is
  finite: $\E_{p_\theta}[\|\partial_i \ell_\theta\|^2] < \infty$ for all
  $\theta \in \Omega^{\mathrm{ideal}}$ and all $i$.  This ensures the Fisher
  metric is well-defined and the Cram\'er--Rao bound applies.
\end{itemize}
\end{mdframed}
\end{samepage}

\paragraph{Relationship to parent assumptions.}
Assumptions IG-A1 through IG-A6 refine the parent framework's assumptions
A1--A6 (see \cite{davis2025manifold}, Section~2.4) to the statistical-manifold
setting.  IG-A1 (statistical structure) is the new condition that enables the
Fisher identification; it is satisfied in every concrete Davis system (HERALD,
VIDAR, AMC, PRISM, KRAKEN).  IG-A2 (Fisher regularity) implies the parent's
smooth chartability condition.  IG-A4 (representation-independence) is the
parent's principle stated in information-geometric language.  IG-A5 (bounded
curvature) specializes the parent's A2 to the Fisher metric.

\subsection{Parameter scales}\label{subsec:param-scales}

Table~\ref{tab:fisher-params} lists illustrative parameter scales for the
Fisher-metric identification in two running examples.

\begin{table}[htbp]
\centering
\caption{\textbf{Illustrative Fisher-metric parameter scales.}  Values are
representative orders of magnitude, not empirical measurements.}
\label{tab:fisher-params}
\begin{tabular}{llcc}
\toprule
Symbol & Meaning & HERALD & VIDAR \\
\midrule
$d$ & Distribution dimension & $\sim 10^3$ (lineage posteriors) & $V - 1$ (vocabulary) \\
$\kappa(\gF)$ & Fisher condition number & $\sim 10$--$50$ & $\sim 5$--$20$ \\
$K_{\max}$ & Curvature bound & $\sim 0.1$--$1.0$ & $\sim 0.01$--$0.1$ \\
$D_{\mathrm{F}}$ & Fisher-Rao diameter & $\sim 2$--$5$ (nats) & $\sim 1$--$3$ (nats) \\
$\mu_1^{\mathrm{F}}$ & Fisher spectral gap & (to be computed) & (to be computed) \\
$\alpha_{\mathrm{eff}}$ & Effective $\alpha$ & $\approx 1$ (exponential) & $\approx 1$ (softmax) \\
\bottomrule
\end{tabular}
\end{table}

%====================================================
\section{Statistical manifolds and the Fisher metric}
\label{sec:statistical-manifolds}
%====================================================

We begin by recalling the foundational objects of information geometry in a
form adapted to the Davis framework.  Throughout this section, all definitions
and theorems through Theorem~\ref{thm:cencov} are classical; we restate them
to fix notation and to make precise how they interact with the Davis
framework's existing structures.

\subsection{Statistical manifolds}

\begin{definition}[Statistical manifold]\label{def:stat-manifold}
Let $\Theta \subseteq \RR^d$ be an open domain and let
$\{p_\theta\}_{\theta \in \Theta}$ be a family of probability distributions on
a measurable space $(\mathcal{X}, \mathcal{A})$ such that:
\begin{enumerate}[label=\textnormal{(\roman*)}, leftmargin=*]
  \item For each $x \in \mathcal{X}$, the map $\theta \mapsto p_\theta(x)$ is
  $C^\infty$ on $\Theta$.
  \item The map $\theta \mapsto p_\theta$ is injective (identifiability).
  \item For each $\theta \in \Theta$, the support of $p_\theta$ does not depend
  on $\theta$ (common support).
\end{enumerate}
The \emph{statistical manifold} is the pair $(\mathcal{S}, \iota)$ where
$\mathcal{S} = \Theta$ as a smooth $d$-dimensional manifold and
$\iota : \mathcal{S} \to \mathcal{P}(\mathcal{X})$ is the parameterization
$\iota(\theta) = p_\theta$.
\end{definition}

\paragraph{Example (exponential families).}
A \emph{$d$-dimensional exponential family} on $\mathcal{X}$ is a statistical
manifold of the form
\begin{equation}\label{eq:exp-family}
p_\theta(x)
= h(x) \exp\!\bigl(\theta^i t_i(x) - \psi(\theta)\bigr),
\end{equation}
where $\theta = (\theta^1, \dots, \theta^d)$ are \emph{natural parameters},
$t = (t_1, \dots, t_d)$ are \emph{sufficient statistics}, $h(x)$ is a base
measure, and the \emph{log-partition function}
\begin{equation}\label{eq:log-partition}
\psi(\theta) = \log \int_{\mathcal{X}} h(x)
\exp\!\bigl(\theta^i t_i(x)\bigr)\, dx
\end{equation}
ensures normalization.  We use the Einstein summation convention throughout.

\paragraph{Example (softmax output distributions).}
At layer~$\ell$ and context position~$t$, a transformer with vocabulary
$\mathcal{V}$ (where $|\mathcal{V}| = V$) produces logits
$z_{\ell,t} \in \RR^V$ and outputs
\[
p_{\ell,t}(v) = \frac{\exp(z_{\ell,t}^v)}{\sum_{w \in \mathcal{V}}
\exp(z_{\ell,t}^w)}
= \exp\!\bigl(z_{\ell,t}^v - \psi(z_{\ell,t})\bigr),
\]
where $\psi(z) = \log \sum_w \exp(z^w)$.  This is a $(V{-}1)$-dimensional
exponential family with natural parameter
$\theta = z_{\ell,t}$ (up to the shift ambiguity $z \mapsto z + c\mathbf{1}$,
which is quotiented out by softmax).

\paragraph{Running examples (statistical structure in Davis systems).}
\begin{itemize}[leftmargin=*]
  \item \textbf{HERALD.}  The HERALD manifold embeds viral spike sequences into
  $\RR^d$ via a pullback metric $g = J_\theta^\top J_\theta$.  Each embedded
  point $z = f_\theta(x)$ parameterizes a posterior distribution over compatible
  viral lineages given the observed antigenic distances.  The statistical manifold
  is the space of these posteriors, with $\mathcal{X}$ the space of
  neutralization assay outcomes and $p_z(\cdot)$ the posterior over lineages
  given embedding~$z$.  Condition~(i) of Definition~\ref{def:stat-manifold}
  holds when the encoder $f_\theta$ is smooth (guaranteed by the Jacobian
  regularizer $\mathcal{L}_{\mathrm{smooth}}$ of Theorem~1 in
  \cite{davis2025manifold}); condition~(ii) holds when the encoding is
  identifiable on the operational region $\Omega_{\mathrm{in}}$.

  \item \textbf{VIDAR.}  Face-recognition embeddings on $\mathbb{S}^{d-1}$
  parameterize von Mises--Fisher (vMF) distributions over identity observations.
  For a unit vector $\mu \in \mathbb{S}^{d-1}$ and concentration $\kappa > 0$,
  the vMF density is
  $p_{\mu,\kappa}(x) = C_d(\kappa) \exp(\kappa \mu^\top x)$,
  which is an exponential family with natural parameter
  $\theta = \kappa\mu$ and sufficient statistic $t(x) = x$.
  The Fisher metric of the vMF family at fixed $\kappa$ is proportional to the
  round metric on $\mathbb{S}^{d-1}$:
  $\gF_{ij}(\mu) = \kappa A_d(\kappa)\,(\delta_{ij} - \mu_i\mu_j)$,
  where $A_d(\kappa) = I_{d/2}(\kappa)/I_{d/2-1}(\kappa)$ is the ratio of
  modified Bessel functions.  Thus VIDAR's inherited spherical geometry is
  already Fisher-canonical for its output family.

  \item \textbf{AMC.}  When lifted via the functor $F_S$, AMC's translator graph
  becomes a manifold whose points parameterize reconstruction-error distributions
  across telemetry formats.  The Fisher metric on this space measures the
  sensitivity of format translation to perturbations in the underlying system
  state---exactly the distortion control that the AMC drift budget $\varepsilon$
  is designed to enforce.
\end{itemize}

\subsection{The Fisher information metric}

\begin{definition}[Fisher information metric]\label{def:fisher-metric}
Let $(\mathcal{S}, \iota)$ be a statistical manifold with log-likelihood
$\ell_\theta(x) := \log p_\theta(x)$.  The \emph{Fisher information metric} is
the symmetric $(0,2)$-tensor field on $\mathcal{S}$ defined by
\begin{equation}\label{eq:fisher-metric}
\gF_{ij}(\theta)
:= \E_{p_\theta}\!\bigl[\partial_i \ell_\theta \cdot \partial_j \ell_\theta\bigr]
= -\E_{p_\theta}\!\bigl[\partial_i \partial_j \ell_\theta\bigr],
\end{equation}
where $\partial_i := \partial / \partial\theta^i$ and the second equality holds
under the common-support and smoothness conditions of
Definition~\ref{def:stat-manifold}.
\end{definition}

The Fisher metric $\gF$ is symmetric, non-negative, and positive definite when
the score functions $\partial_i \ell_\theta$ are linearly independent---which
holds generically on identifiable families.  It thus equips $\mathcal{S}$ with
a Riemannian structure.

\begin{proposition}[Fisher metric on exponential families]
\label{prop:fisher-exp-family}
On an exponential family~\eqref{eq:exp-family} with log-partition function
$\psi(\theta)$, the Fisher metric is the Hessian of $\psi$:
\begin{equation}\label{eq:fisher-hessian}
\gF_{ij}(\theta) = \frac{\partial^2 \psi}{\partial\theta^i \partial\theta^j}
= \Cov_{p_\theta}\!\bigl[t_i, t_j\bigr].
\end{equation}
In particular, $\gF$ is the covariance matrix of the sufficient statistics
under~$p_\theta$.
\end{proposition}

\begin{proof}
For an exponential family,
$\ell_\theta(x) = \log h(x) + \theta^i t_i(x) - \psi(\theta)$, so
$\partial_i \ell_\theta = t_i(x) - \partial_i \psi(\theta)$.  Then
$\gF_{ij} = \E[(\partial_i \ell_\theta)(\partial_j \ell_\theta)]
= \E[(t_i - \E[t_i])(t_j - \E[t_j])]
= \Cov[t_i, t_j]$.  Since
$\E[t_i] = \partial_i \psi$ and
$\Cov[t_i, t_j] = \partial_i \partial_j \psi$, the result follows.
\end{proof}

\paragraph{Explicit computation: softmax Fisher metric.}
For the softmax exponential family on vocabulary $\mathcal{V}$ with
$|\mathcal{V}| = V$ and natural parameters $\theta = z$ (logits), the
log-partition function is $\psi(\theta) = \log \sum_{v=1}^V \exp(\theta^v)$.
By Proposition~\ref{prop:fisher-exp-family},
\begin{equation}\label{eq:softmax-fisher-explicit}
\gF_{ij}(z) = \frac{\partial^2 \psi}{\partial z^i \partial z^j}
= \delta_{ij}\, p_i - p_i\, p_j
= \begin{cases}
p_i(1-p_i) & \text{if } i = j, \\
-p_i\, p_j & \text{if } i \neq j,
\end{cases}
\end{equation}
where $p_i = \exp(z^i)/\sum_w \exp(z^w)$.  In matrix form,
$\gF(z) = \diag(p) - p\, p^\top$.

This is the covariance matrix of the categorical distribution and can be
computed \emph{exactly} from the softmax output---no stochastic estimation,
no sampling, no approximation.  For a vocabulary of size $V = 50{,}000$, the
Fisher matrix is $50{,}000 \times 50{,}000$ but has rank at most $V-1$ and
admits efficient storage and manipulation via the Sherman--Morrison structure
$\gF = \diag(p) - p\,p^\top$.

\paragraph{Numerical example (toy vocabulary).}
Consider a toy model with $V = 4$ and logits $z = (2.0, 1.0, 0.5, 0.0)$.  The
softmax probabilities are $p \approx (0.506, 0.186, 0.113, 0.068)$.  The
Fisher matrix is
\[
\gF \approx
\begin{pmatrix}
0.250 & -0.094 & -0.057 & -0.035 \\
-0.094 & 0.152 & -0.021 & -0.013 \\
-0.057 & -0.021 & 0.100 & -0.008 \\
-0.035 & -0.013 & -0.008 & 0.064
\end{pmatrix}.
\]
The eigenvalues are approximately $(0.341, 0.157, 0.068, 0.000)$---the zero
eigenvalue corresponds to the shift ambiguity $z \mapsto z + c\mathbf{1}$,
confirming the $(V{-}1)$-dimensional structure.  The condition number on the
non-degenerate subspace is $\kappa \approx 0.341/0.068 \approx 5.0$, which is
moderate and well within the Fisher regularity assumption IG-A2.

\paragraph{Fisher metric of the vMF family (VIDAR).}
For a von Mises--Fisher distribution on $\mathbb{S}^{d-1}$ with concentration
$\kappa$ and mean direction $\mu$, the Fisher metric restricted to the
tangent space $T_\mu \mathbb{S}^{d-1}$ at fixed $\kappa$ is
\begin{equation}\label{eq:vmf-fisher}
\gF_{ij}(\mu) = \kappa A_d(\kappa)\,(\delta_{ij} - \mu_i \mu_j),
\end{equation}
where $A_d(\kappa) = I_{d/2}(\kappa)/I_{d/2-1}(\kappa)$.  The projector
$\Pi_\mu = I - \mu\mu^\top$ restricts to the tangent space, so this is simply
$\kappa A_d(\kappa)$ times the round metric on the sphere.  For large $\kappa$
(high concentration), $A_d(\kappa) \to 1$ and the Fisher metric approaches
$\kappa$ times the round metric.

This confirms the observation from Section~\ref{subsec:running-examples}: VIDAR's
inherited spherical geometry is Fisher-canonical for its output family.

\subsection{\v{C}encov's uniqueness theorem and the core identification}

The Fisher metric is not one of many possible choices.  It is the
\emph{only} Riemannian metric on statistical manifolds that respects the
structure of statistical inference.

\begin{theorem}[\v{C}encov uniqueness, 1982]\label{thm:cencov}
On the category of finite statistical models with morphisms given by
congruent Markov kernels (sufficient statistic reductions), the Fisher
information metric $\gF$ is the unique Riemannian metric, up to a positive
scalar multiple, that is invariant under all such morphisms.
\end{theorem}

\v{C}encov's original proof~\cite{cencov1982} established uniqueness on
finite sample spaces.  Extensions to general dominated families were given
by Campbell~\cite{campbell1986} and Ay et al.~\cite{ay2017}; the essential
content is unchanged: any Riemannian metric satisfying statistical invariance
must be proportional to the Fisher metric.

We now state the core identification of this branch.

\begin{mdframed}[backgroundcolor=gray!10, linewidth=0.4pt]
\noindent\textbf{Core Identification (informal preview).}
\emph{The Davis framework's representation-independence principle requires the
metric $g$ to be invariant under changes of feature-space coordinates that
preserve sufficient statistics of the sameness structure~$S$.  \v{C}encov's
uniqueness theorem requires the same invariance and concludes that the only
such metric is the Fisher information metric $\gF$, up to a positive scalar.
These are the same requirement.  Therefore $g = c \cdot \gF$ is forced, not
chosen.}
\end{mdframed}
\vspace{1ex}

\begin{corollary}[Unique Davis metric in the statistical-manifold regime]\label{cor:canonical-davis}
Let $(M, g)$ be a Davis manifold whose points parameterize probability
distributions (i.e., $M$ carries the structure of a statistical manifold
$\mathcal{S}$).  If $g$ is required to be invariant under sufficient statistic
reductions of the observation model---the invariance demanded by the Davis
framework's representation-independence principle---then
\[
g = c \cdot \gF
\]
for some constant $c > 0$.
\end{corollary}

\begin{proof}
The Davis framework's representation-independence principle requires that
detection guarantees be invariant under reparameterizations that preserve
sufficient statistics~\cite{davis2025manifold, davis2025sameness}.  A
Riemannian metric is invariant under such reparameterizations if and only if it
is invariant under congruent Markov kernels on the induced statistical model.
By Theorem~\ref{thm:cencov}, the only such metric is $c \cdot \gF$ for
$c > 0$.
\end{proof}

\begin{remark}[What this identification is and is not]
\label{rmk:identification}
Corollary~\ref{cor:canonical-davis} is not a deep new result.  It is a
recognition that two independently motivated invariance conditions---one from
detection theory (Davis representation-independence), one from mathematical
statistics (\v{C}encov invariance under sufficient statistics)---impose the
same constraint on the metric.  The contribution is not the mathematics but the
\emph{assembly}: placing the Fisher metric into the Davis framework reveals
that the parent's curvature, holonomy, spectral gap, and capacity bounds all
have canonical statistical interpretations.
\end{remark}

%====================================================
\section{The \texorpdfstring{$\alpha$}{alpha}-connection family}
\label{sec:alpha-connections}
%====================================================

The Fisher metric alone recovers the parent framework's metric.  But
statistical manifolds carry additional structure: a one-parameter family of
affine connections, introduced by Amari~\cite{amari1985}, that includes the
Levi-Civita connection as a special case.  This family has no analog in the
parent framework and constitutes the genuine structural extension of Branch~X.

\subsection{Definition and duality}

\begin{definition}[$\alpha$-connections]\label{def:alpha-connection}
Let $(\mathcal{S}, \gF)$ be a statistical manifold.  For $\alpha \in \RR$, the
\emph{$\alpha$-connection} $\nabla^{(\alpha)}$ is the affine connection on
$\mathcal{S}$ defined by
\begin{equation}\label{eq:alpha-connection}
\gF\!\left(\nabla^{(\alpha)}_{\partial_i} \partial_j,\; \partial_k\right)
= \E_{p_\theta}\!\left[\left(\partial_i \partial_j \ell_\theta
+ \frac{1-\alpha}{2}\, \partial_i \ell_\theta \cdot
\partial_j \ell_\theta\right) \partial_k \ell_\theta\right].
\end{equation}
\end{definition}

Three values are distinguished:
\begin{itemize}
  \item $\alpha = 0$: the \emph{Levi-Civita connection} $\nabla^{(0)}$, which
  is the unique torsion-free, metric-compatible connection on
  $(\mathcal{S}, \gF)$.  This is the connection used throughout the parent
  framework.
  \item $\alpha = +1$: the \emph{exponential connection} $\nabla^{(e)}$, which
  is flat on exponential families.
  \item $\alpha = -1$: the \emph{mixture connection} $\nabla^{(m)}$, which is
  flat on mixture families.
\end{itemize}

The central structural result is duality.

\begin{theorem}[Conjugate duality]\label{thm:duality}
For every $\alpha \in \RR$, the pair
$(\nabla^{(\alpha)}, \nabla^{(-\alpha)})$ is conjugate with respect to the
Fisher metric: for any smooth vector fields $X, Y, Z$ on $\mathcal{S}$,
\begin{equation}\label{eq:duality}
X \cdot \gF(Y, Z)
= \gF\!\left(\nabla^{(\alpha)}_X Y,\; Z\right)
+ \gF\!\left(Y,\; \nabla^{(-\alpha)}_X Z\right).
\end{equation}
In particular, $\nabla^{(0)}$ is self-dual, and the pair
$(\nabla^{(e)}, \nabla^{(m)})$ is dual.
\end{theorem}

\begin{proof}
Direct computation from Definition~\ref{def:alpha-connection}.  Write
$\Gamma^{(\alpha)}_{ij,k}$ for the Christoffel symbols of
$\nabla^{(\alpha)}$ lowered by $\gF$.  Then
\begin{align*}
\Gamma^{(\alpha)}_{ij,k} + \Gamma^{(-\alpha)}_{ik,j}
&= \E\!\left[\left(\partial_i \partial_j \ell + \frac{1-\alpha}{2}
\partial_i \ell\, \partial_j \ell\right) \partial_k \ell\right] \\
&\quad + \E\!\left[\left(\partial_i \partial_k \ell + \frac{1+\alpha}{2}
\partial_i \ell\, \partial_k \ell\right) \partial_j \ell\right].
\end{align*}
Adding these and using
$\partial_i(\gF_{jk}) = \E[\partial_i \partial_j \ell\, \partial_k \ell]
+ \E[\partial_j \ell\, \partial_i \partial_k \ell]
+ \E[\partial_i \ell\, \partial_j \ell\, \partial_k \ell]$
(which holds by differentiating under the expectation and using the score
equation $\E[\partial_i \ell] = 0$), one verifies that
$\Gamma^{(\alpha)}_{ij,k} + \Gamma^{(-\alpha)}_{ik,j}
= \partial_i(\gF_{jk})$,
which is exactly the duality condition~\eqref{eq:duality}.
\end{proof}

\noindent\emph{Intuition.}  Duality means that the two connections
$\nabla^{(\alpha)}$ and $\nabla^{(-\alpha)}$ ``share'' the metric: they are
compatible with $\gF$ not individually (neither is metric-compatible unless
$\alpha = 0$) but \emph{jointly}: their failure to preserve the metric cancels
when combined.  For a vector $Y$ parallel-transported along $X$ using
$\nabla^{(\alpha)}$, the inner product $\gF(Y, Z)$ changes; but if $Z$ is
simultaneously transported using $\nabla^{(-\alpha)}$, the changes cancel and
the inner product is preserved.

This is the geometric content of conjugate pairs in estimation theory: the
exponential family's natural parameterization and mean parameterization are
``opposite'' in a precise sense, and neither alone preserves the Fisher inner
product, but together they do.

\subsection{The Amari--Chentsov tensor}

The entire $\alpha$-connection family is determined by the Fisher metric $\gF$
and a single additional object: the \emph{skewness tensor} (or
Amari--Chentsov tensor).

\begin{definition}[Amari--Chentsov tensor]\label{def:ac-tensor}
The \emph{Amari--Chentsov tensor} is the symmetric $(0,3)$-tensor field
$T_{ijk}$ on $(\mathcal{S}, \gF)$ defined by
\begin{equation}\label{eq:ac-tensor}
T_{ijk}(\theta) := \E_{p_\theta}\!\left[
\partial_i \ell_\theta \cdot \partial_j \ell_\theta \cdot
\partial_k \ell_\theta\right].
\end{equation}
This is the third central moment of the score function.
\end{definition}

\begin{lemma}[$\alpha$-connections from the Amari--Chentsov tensor]
\label{lem:alpha-from-ac}
The Christoffel symbols of $\nabla^{(\alpha)}$ (lowered by $\gF$) are related
to those of the Levi-Civita connection $\nabla^{(0)}$ by
\begin{equation}\label{eq:alpha-from-ac}
\Gamma^{(\alpha)}_{ij,k}
= \Gamma^{(0)}_{ij,k} - \frac{\alpha}{2}\, T_{ijk}.
\end{equation}
\end{lemma}

\begin{proof}
Direct computation from Definition~\ref{def:alpha-connection}.  The
Levi-Civita Christoffel symbols satisfy
$\Gamma^{(0)}_{ij,k} = \frac{1}{2}(\partial_i \gF_{jk} + \partial_j \gF_{ik}
- \partial_k \gF_{ij})$.  The difference
$\Gamma^{(\alpha)}_{ij,k} - \Gamma^{(0)}_{ij,k}$ is computed from
\eqref{eq:alpha-connection} by expanding the expectation and using the score
equation $\E[\partial_i \ell] = 0$.  The result is
$-\frac{\alpha}{2} T_{ijk}$.
\end{proof}

\noindent\emph{Significance.}  Lemma~\ref{lem:alpha-from-ac} shows that the
\emph{entire} $\alpha$-connection family is generated by two objects:
$\gF$ and $T$.  The Fisher metric is the ``metric part'' and the
Amari--Chentsov tensor is the ``cubic part.''  On exponential families,
$T_{ijk}$ has a simple closed form in natural coordinates:
$T_{ijk} = -\partial_i \partial_j \partial_k \psi(\theta)$, the third
derivative of the log-partition function.  When $T \equiv 0$ (Gaussian
families in natural parameterization), all $\alpha$-connections coincide with
the Levi-Civita connection and the manifold is ``statistically flat.''

\begin{remark}[What duality means for the Davis framework]
\label{rmk:duality-davis}
The parent framework's curvature, parallel transport, geodesics, and holonomy
are all defined using the Levi-Civita connection $\nabla^{(0)}$.  Duality
reveals that $\nabla^{(0)}$ is the ``average'' of two flat connections
$\nabla^{(e)}$ and $\nabla^{(-e)}$ on exponential families.  The parent's
curvature is entirely an artifact of this averaging: it measures the cost of
using a single symmetric connection rather than exploiting the exponential
family structure.  On a non-exponential manifold, all three connections are
curved, and the $\alpha$-curvature family provides a finer diagnostic than the
Levi-Civita curvature alone.
\end{remark}

\subsection{Dual coordinate systems on exponential families}

On exponential families, the dual connections induce canonical coordinate
systems via the Legendre transform.

\begin{theorem}[Dual coordinates]\label{thm:dual-coordinates}
Let $\mathcal{S}$ be an exponential family~\eqref{eq:exp-family} with
log-partition function $\psi(\theta)$.  Define:
\begin{enumerate}[label=\textnormal{(\roman*)}, leftmargin=*]
  \item \emph{Natural parameters} (or $\theta$-coordinates):
  the coordinates in which $\nabla^{(e)}$ is flat.
  \item \emph{Expectation parameters} (or $\eta$-coordinates):
  \begin{equation}\label{eq:eta-coords}
  \eta^i := \E_{p_\theta}[t_i]
  = \frac{\partial \psi}{\partial \theta^i},
  \end{equation}
  which are the coordinates in which $\nabla^{(m)}$ is flat.
\end{enumerate}
The two coordinate systems are related by the Legendre transform:
\begin{equation}\label{eq:legendre}
\varphi(\eta) := \sup_\theta\!\bigl\{\theta \cdot \eta - \psi(\theta)\bigr\}
= \theta(\eta) \cdot \eta - \psi\bigl(\theta(\eta)\bigr),
\end{equation}
where $\varphi(\eta)$ is the negative entropy (or dual potential).  The
Fisher metric in the two coordinate systems is:
\begin{equation}\label{eq:fisher-dual}
\gF_{ij}(\theta) = \frac{\partial^2 \psi}{\partial\theta^i\, \partial\theta^j},
\qquad
g^{\mathrm{F},ij}(\eta) = \frac{\partial^2 \varphi}{\partial\eta^i\, \partial\eta^j}.
\end{equation}
\end{theorem}

\begin{proof}
That $\eta^i = \partial_i \psi$ follows from
$\E[t_i] = \partial_i \log Z = \partial_i \psi$.  The Legendre duality between
$\psi$ and $\varphi$ is standard convex analysis: $\psi$ is strictly convex
on the natural parameter space (since $\gF = \nabla^2 \psi > 0$), so the
Legendre transform is well-defined and involutive.  The Fisher metric in
$\eta$-coordinates is $g^{\mathrm{F},ij} = \partial^2 \varphi / \partial\eta^i
\partial\eta^j$ by the standard identity for Legendre-dual Hessians.
\end{proof}

\paragraph{Worked example: dual coordinates on softmax.}
For the softmax family with $V = 4$ and natural parameters
$\theta = (z^1, z^2, z^3, z^4)$:
\begin{itemize}[leftmargin=*]
  \item \emph{Natural coordinates} ($\theta$-coordinates):
  the logits $z = (2.0, 1.0, 0.5, 0.0)$.  The exponential connection
  $\nabla^{(e)}$ is flat in these coordinates: Christoffel symbols vanish.
  Straight lines in logit space are $e$-geodesics.

  \item \emph{Expectation coordinates} ($\eta$-coordinates):
  the softmax probabilities $\eta = p = (0.506, 0.186, 0.113, 0.068)$.  The
  mixture connection $\nabla^{(m)}$ is flat in these coordinates: straight
  lines in probability space are $m$-geodesics (i.e., mixture interpolations
  $q_t = (1-t)p + tq$ are $m$-geodesics).

  \item \emph{Log-partition function}:
  $\psi(\theta) = \log(\exp(2) + \exp(1) + \exp(0.5) + 1) \approx 2.695$.

  \item \emph{Dual potential}:
  $\varphi(\eta) = \sum_v \eta^v \log \eta^v \approx -1.165$ (the negative
  entropy of the softmax distribution).

  \item \emph{Legendre relation}:
  $\theta \cdot \eta - \psi(\theta) = (2)(0.506) + (1)(0.186) + (0.5)(0.113) +
  (0)(0.068) - 2.695 \approx -1.165 = \varphi(\eta)$. \checkmark
\end{itemize}

\paragraph{What dual coordinates mean for Davis systems.}
In a transformer, the logits are natural coordinates and the softmax outputs
are expectation coordinates.  The softmax function itself is the
Legendre-transform map $\theta \mapsto \eta = \nabla \psi(\theta)$.  This
means the transformer's computation naturally alternates between two
``flat'' coordinate systems:
\begin{itemize}[leftmargin=*]
  \item Logit space (before softmax): flat under $\nabla^{(e)}$, natural for
  exponential-family operations (inner products, additions in log-space).
  \item Probability space (after softmax): flat under $\nabla^{(m)}$, natural
  for mixture operations (convex combinations, averaging).
\end{itemize}
The Levi-Civita connection $\nabla^{(0)}$---the parent framework's
connection---is curved in both coordinate systems.  This curvature is the
``price of symmetry'': using a single self-dual connection that treats both
parameterizations equally.

\begin{remark}[Bridge to Branch~VI]\label{rmk:bridge-vi}
The log-partition function $\psi(\theta) = \log Z(\theta)$ is the
\emph{free energy} (up to a temperature factor) in the statistical-mechanics
formulation of Branch~VI.  Proposition~\ref{prop:fisher-exp-family} and
Theorem~\ref{thm:dual-coordinates} together say: the Fisher metric is the
Hessian of the free energy, i.e., the \emph{susceptibility matrix} of the
thermal ensemble.  This connection is structural, not analogical: the same
function $\psi$ governs both the information geometry and the thermodynamics.
\end{remark}

\subsection{Flatness and curvature of \texorpdfstring{$\alpha$}{alpha}-connections}

\begin{theorem}[Exponential flatness]\label{thm:e-flat}
On an exponential family $\mathcal{S}$, the exponential connection
$\nabla^{(e)}$ is flat: its Riemann curvature tensor vanishes identically,
$R^{(e)} \equiv 0$.  In natural coordinates $\theta$, the Christoffel symbols
of $\nabla^{(e)}$ vanish:
$\Gamma^{(e)\, k}_{ij}(\theta) = 0$ for all $i, j, k$.
\end{theorem}

\begin{proof}
In natural coordinates, $\ell_\theta(x) = \theta^i t_i(x) - \psi(\theta)
+ \log h(x)$, so $\partial_i \partial_j \ell_\theta = -\partial_i \partial_j
\psi = -\gF_{ij}$.  The $\alpha = 1$ Christoffel symbols
(lowered) are
$\Gamma^{(e)}_{ij,k} = \E[(\partial_i \partial_j \ell) \partial_k \ell]
= -\gF_{ij} \cdot \E[\partial_k \ell] = 0$
since $\E[\partial_k \ell] = 0$.  Therefore $\nabla^{(e)}$ is flat in natural
coordinates, and flatness is coordinate-independent.
\end{proof}

\begin{corollary}[Mixture flatness]\label{cor:m-flat}
By duality (Theorem~\ref{thm:duality}), the mixture connection $\nabla^{(m)}$
is flat on exponential families when expressed in expectation coordinates
$\eta$.
\end{corollary}

The Levi-Civita connection $\nabla^{(0)}$, by contrast, is generically
\emph{curved} on exponential families.  Its curvature is the ``price of
symmetry'': the cost of insisting on a single self-dual connection rather than
choosing one of the two flat ones.

%====================================================
\section{Curvature, capacity, and the Cram\'er--Rao connection}
\label{sec:curvature-capacity}
%====================================================

We now show that the parent framework's curvature capacity law is a
consequence of the Cram\'er--Rao inequality on the Fisher manifold, and that
exponential families saturate the bound.

\subsection{The Cram\'er--Rao inequality on statistical manifolds}

\begin{theorem}[Cram\'er--Rao bound]\label{thm:cramer-rao}
Let $(\mathcal{S}, \gF)$ be a statistical manifold and let
$\hat\theta : \mathcal{X} \to \RR^d$ be an unbiased estimator of $\theta$.
Then the covariance matrix of $\hat\theta$ satisfies
\begin{equation}\label{eq:cramer-rao}
\Cov_{p_\theta}[\hat\theta]
\;\succeq\;
\bigl(\gF(\theta)\bigr)^{-1}
\end{equation}
in the Loewner (positive semidefinite) order.  Equality holds if and only if
$\hat\theta$ is a sufficient statistic and $\mathcal{S}$ is an exponential
family.
\end{theorem}

\begin{proof}
Classical; see e.g.\ Rao (1945) or Lehmann \& Casella (1998).
\end{proof}

Efron~\cite{efron1975} showed that on curved statistical manifolds (where
$R^{(e)} \neq 0$), the Cram\'er--Rao bound tightens by a curvature correction:
the variance of any estimator exceeds $(\gF)^{-1}$ by an amount that depends on
the \emph{statistical curvature} of the model.  We use this to derive the
curvature capacity bound.

\begin{lemma}[Efron curvature correction]\label{lem:efron}
Let $(\mathcal{S}, \gF)$ be a $d$-dimensional statistical manifold and let
$\hat\theta$ be an asymptotically efficient estimator.  The covariance of
$\hat\theta$ satisfies
\begin{equation}\label{eq:efron-correction}
\Cov[\hat\theta] = (\gF)^{-1}
+ \frac{1}{n}\,(\gF)^{-1}\, C_E\, (\gF)^{-1}
+ O(n^{-2}),
\end{equation}
where $n$ is the sample size and $C_E$ is Efron's curvature matrix:
\[
(C_E)_{ij} = \sum_{k,\ell}
(\gF)^{k\ell}\, R^{(e)}_{ik\ell j},
\]
with $R^{(e)}$ the Riemann curvature tensor of the exponential connection
$\nabla^{(e)}$.  On exponential families, $R^{(e)} \equiv 0$ (by
Theorem~\ref{thm:e-flat}), so $C_E = 0$ and the Cram\'er--Rao bound is
achieved to second order.
\end{lemma}

\begin{proof}[Proof sketch]
Efron's original argument~\cite{efron1975} uses the geometry of the score
function manifold: the locus $\{s_\theta := \nabla \ell_\theta\}$ in
$L^2(p_\theta)$ forms a surface whose curvature determines the gap between
the Cram\'er--Rao bound and the achievable variance.  The curvature of this
surface is exactly the curvature of $\nabla^{(e)}$ on $\mathcal{S}$.  When
$\nabla^{(e)}$ is flat (exponential families), the score surface is affinely
embedded and the bound is tight.  A complete treatment is given in
Amari~\cite{amari1985}, Chapter~4.
\end{proof}

\subsection{The curvature capacity bound}

\begin{theorem}[Curvature capacity from Cram\'er--Rao]
\label{thm:curvature-capacity}
Let $(M, \gF)$ be a Davis manifold identified with a $d$-dimensional
statistical manifold $(\mathcal{S}, \gF)$ via
Corollary~\ref{cor:canonical-davis}.  Suppose
$K^{(0)}(\sigma) \ge \kappa > 0$
for all tangent $2$-planes $\sigma$ on $\mathcal{S}$ (i.e., the sectional
curvature is bounded below by $\kappa$).

\begin{enumerate}[label=\textnormal{(\roman*)}, leftmargin=*]
  \item \textbf{Path capacity.}  For any completion path
  $\gamma : [0,1] \to M$, the
  $\delta$-packing number of $\gamma$ in the Fisher-Rao metric satisfies
  \begin{equation}\label{eq:capacity-path}
  C_\delta(\gamma) \;\le\;
  \frac{\pi}{\delta\,\sqrt{\kappa}} + 1.
  \end{equation}

  \item \textbf{Region capacity (Davis curvature capacity law).}
  For any completion region $\Omega \subseteq \mathcal{S}$, the
  $\delta$-packing number satisfies
  \begin{equation}\label{eq:capacity-region}
  C_\delta(\Omega) \;\le\;
  c(d)\;\kappa^{-d/2}\;\delta^{-d},
  \end{equation}
  where $c(d) = 2^{d+1}\pi^{(d+1)/2} /
  (\Gamma(\frac{d+1}{2})\,\omega_d)$ is an explicit dimensional constant
  and $\omega_d = \pi^{d/2}/\Gamma(d/2+1)$ is the Euclidean unit-ball
  volume.  In particular, for $d = 2$ this contains the parent framework's
  curvature capacity law $C \lesssim 1/K$ as a specialization:
  \begin{equation}\label{eq:capacity-law}
  C_\delta(\Omega) \;\le\;
  c(2)\;\frac{1}{\kappa\,\delta^2}.
  \end{equation}

  \item \textbf{Efron correction on curved models.}  If the statistical
  manifold has nonzero Efron curvature $c_E$ (i.e., the model is
  not an exponential family), the effective resolution enlarges to
  $\delta_{\mathrm{eff}} = \delta\,\sqrt{1 + \|c_E\|^2}$, and both
  bounds above hold with $\delta$ replaced by $\delta_{\mathrm{eff}}$:
  \begin{equation}\label{eq:capacity-efron}
  C_{\delta}(\Omega)
  \;\le\;
  c(d)\;\frac{1}
  {\kappa^{d/2}\;\delta^d\;
  (1 + \sup_{\theta \in \Omega} \|c_E(\theta)\|^2)^{d/2}}.
  \end{equation}
\end{enumerate}

In all cases, higher Levi-Civita curvature reduces capacity.  At general
$d$, the bound $C \sim \kappa^{-d/2}$ is \emph{stronger} than the parent's
$1/K$ law.
\end{theorem}

\begin{proof}
We prove each part.

\paragraph{Step 1: Diameter bound (Bonnet--Myers).}
Since $K^{(0)}(\sigma) \ge \kappa > 0$ for all tangent $2$-planes,
the Ricci curvature satisfies
$\mathrm{Ric}_{\gF} \ge (d{-}1)\kappa > 0$.
By the Bonnet--Myers theorem~\cite{bonnet1855, myers1941},
$\mathcal{S}$ is compact with
\begin{equation}\label{eq:bonnet-myers}
\mathrm{diam}(\mathcal{S}, \gF) \le \frac{\pi}{\sqrt{\kappa}}.
\end{equation}

\paragraph{Step 2: Path packing (part~(i)).}
Points along a curve of Fisher-Rao arc length
$\ell_{\mathrm{F}}(\gamma) \le \pi/\sqrt\kappa$ that are mutually
$\delta$-separated number at most
$\lfloor \ell_{\mathrm{F}} / \delta \rfloor + 1
\le \pi/(\delta\sqrt\kappa) + 1$.  This gives~\eqref{eq:capacity-path}.

\paragraph{Step 3: Volume packing (part~(ii)).}
In a $\delta$-packing of $\Omega$, the balls $B(x_i, \delta/2)$ are
disjoint and contained in $\mathcal{S}$, so
\[
C_\delta(\Omega)
\;\le\;
\frac{\mathrm{Vol}_{\gF}(\mathcal{S})}
     {\inf_{x \in \Omega}\mathrm{Vol}_{\gF}(B(x,\,\delta/2))}.
\]
By the Bishop--Gromov volume comparison theorem, on a $d$-manifold with
$K^{(0)} \ge \kappa > 0$:
\begin{itemize}[leftmargin=*]
\item The total volume is bounded above by that of the round
$d$-sphere $S^d(1/\sqrt\kappa)$:
\[
\mathrm{Vol}_{\gF}(\mathcal{S})
\;\le\;
\mathrm{Vol}(S^d(1/\sqrt\kappa))
\;=\;
\frac{2\,\pi^{(d+1)/2}}
     {\Gamma\!\bigl(\tfrac{d+1}{2}\bigr)}\;
\kappa^{-d/2}.
\]
\item For $\delta/2$ below the injectivity radius
($\le \pi/\sqrt\kappa$ by Klingenberg), the small-ball volume
expansion gives
\[
\mathrm{Vol}_{\gF}(B(x, \delta/2))
\;\ge\;
\omega_d\;\Bigl(\frac{\delta}{2}\Bigr)^d
\;\Bigl(1 - \frac{d\,\kappa\,\delta^2}{24(d+2)}
+ O(\delta^4)\Bigr),
\]
where $\omega_d = \pi^{d/2}/\Gamma(d/2+1)$.
\end{itemize}
Combining and dropping the lower-order correction:
\[
C_\delta(\Omega)
\;\le\;
\frac{2^{d+1}\,\pi^{(d+1)/2}}
     {\Gamma\!\bigl(\tfrac{d+1}{2}\bigr)\;\omega_d}
\;\kappa^{-d/2}\;\delta^{-d}
\;=\; c(d)\;\kappa^{-d/2}\;\delta^{-d}.
\]
At $d = 2$: $c(2) = 2^3 \pi^{3/2}/(\Gamma(3/2)\,\pi) = 8$, and
$\kappa^{-d/2} = \kappa^{-1}$, giving the parent's
$C \le 8\,K^{-1}\,\delta^{-2}$.

\paragraph{Step 4: Efron correction (part~(iii)).}
By Efron~\cite{efron1975}, on a non-exponential statistical manifold with
Efron curvature $c_E(\theta) := \|R^{(e)}(\theta)\|$, the minimum variance
of any unbiased estimator exceeds
$(\gF)^{-1}$ by a multiplicative factor $(1 + \|c_E\|^2)$.
This enlarges the minimum $\delta$-distinguishability separation to
$\delta_{\mathrm{eff}} = \delta\,\sqrt{1 + \|c_E\|^2}$.  Substituting
$\delta_{\mathrm{eff}}$ for $\delta$ in~\eqref{eq:capacity-path}
and~\eqref{eq:capacity-region} gives~\eqref{eq:capacity-efron}.
\end{proof}

\begin{corollary}[Saturation on exponential families]
\label{cor:saturation}
On an exponential family $\mathcal{S}$, the exponential connection is flat
($R^{(e)} = 0$), so $c_E \equiv 0$.  The Efron correction
in~\eqref{eq:capacity-efron} vanishes, and the capacity
bounds~\eqref{eq:capacity-path}--\eqref{eq:capacity-region} are tight: the
maximum likelihood estimator achieves the Cram\'er--Rao bound, and the
packing arguments reach maximum density.\end{corollary}

\begin{proof}
By Theorem~\ref{thm:e-flat}, $R^{(e)} \equiv 0$ on exponential families.
The Efron curvature correction is $\|c_E\| = 0$, so
$\delta_{\mathrm{eff}} = \delta$.
The path bound~\eqref{eq:capacity-path}
is achieved with equality (up to the $+1$ floor term) when $\gamma$ is a
geodesic of length $\pi/\sqrt{\kappa}$.
The region bound~\eqref{eq:capacity-region} is achieved when
$\Omega = \mathcal{S}$ is the full round sphere $S^{V-1}$
(in the square-root parametrization $\xi_i = 2\sqrt{p_i}$), for which
Bishop--Gromov comparison holds with equality.
\end{proof}


\begin{remark}[Why softmax is special]\label{rmk:softmax}
The softmax output layer of a transformer produces an exponential family
distribution~\eqref{eq:exp-family} with natural parameters given by the logits.
Corollary~\ref{cor:saturation} provides a structural explanation for why
softmax-based architectures are \emph{locally} geometrically efficient: at the
output layer, their distributions lie on an $e$-flat manifold where the
capacity bound is tight.  The full transformer (including attention
nonlinearities, layer normalization, and residual connections) is not globally
an exponential family; the saturation claim applies to the
\emph{output distribution conditional on the hidden state}.
Non-exponential output distributions (e.g., those produced by sparse attention
or mixture-of-experts layers) would have nonzero statistical curvature and
correspondingly reduced capacity.
\end{remark}

\subsection{KL divergence as tolerance budget}

\begin{theorem}[KL divergence as infinitesimal Fisher-Rao distance]
\label{thm:kl-fisher}
For nearby distributions $p_\theta$ and $p_{\theta + d\theta}$ on a
statistical manifold $(\mathcal{S}, \gF)$,
\begin{equation}\label{eq:kl-fisher}
\DKL(p_\theta \| p_{\theta + d\theta})
= \tfrac{1}{2}\, \gF_{ij}(\theta)\, d\theta^i\, d\theta^j
+ O(\|d\theta\|^3).
\end{equation}
The KL divergence is the local quadratic approximation to the squared
Fisher-Rao geodesic distance.
\end{theorem}

\begin{proof}
Taylor expand $\DKL(p_\theta \| p_{\theta + d\theta})
= \E_{p_\theta}[\ell_\theta - \ell_{\theta + d\theta}]$
to second order in $d\theta$.  The first-order term vanishes by
$\E[\partial_i \ell_\theta] = 0$, and the second-order term is
$\tfrac{1}{2} \gF_{ij}\, d\theta^i\, d\theta^j$.
\end{proof}

\begin{corollary}[Tolerance budget as KL budget]\label{cor:tolerance-kl}
The parent framework's tolerance budget $\tau$ has a canonical
information-geometric interpretation: it is a KL divergence budget.  The
total information-theoretic distance a completion is allowed to travel is
\[
\sum_{\ell} \DKL(p_{\ell} \| p_{\ell+1}) \;\le\; \tau,
\]
where the sum is over layers (or time steps) of the completion path.
\end{corollary}

\begin{corollary}[Holonomy as information loss]\label{cor:holonomy-kl}
The parent framework's holonomy $\mathrm{Hol}(\gamma)$ around a closed path
$\gamma$ in $(\mathcal{S}, \gF)$ satisfies
\begin{equation}\label{eq:holonomy-area}
\norm{\mathrm{Hol}(\gamma) - I}
\;\le\;
C_R \cdot \mathrm{Area}_{\gF}(\gamma) + O(\mathrm{Area}^2),
\end{equation}
where $C_R$ depends on the Riemann curvature of $\nabla^{(0)}$ and
$\mathrm{Area}_{\gF}$ is the Fisher-metric area enclosed by $\gamma$.  This
is the Ambrose--Singer theorem applied to the Fisher Levi-Civita connection:
information loss around constraint loops is controlled by curvature times
enclosed area.
\end{corollary}


%====================================================
\section{Fisher spectral capacity and the central conjecture}
\label{sec:spectral-capacity}
%====================================================

Branch~IX established that spectral gap controls capacity on CSP graphs.
We now lift this to the Fisher-metric setting.

\subsection{Fisher Laplacian}

\begin{definition}[Fisher Laplacian]\label{def:fisher-laplacian}
On a statistical manifold $(\mathcal{S}, \gF)$, the \emph{Fisher Laplacian}
$\Delta^{\mathrm{F}}$ is the Laplace--Beltrami operator of the Fisher metric:
for $f \in C^\infty(\mathcal{S})$,
\begin{equation}\label{eq:fisher-laplacian}
\Delta^{\mathrm{F}} f
= \frac{1}{\sqrt{\det \gF}}
\partial_i\!\left(\sqrt{\det \gF}\; (\gF)^{ij}\, \partial_j f\right).
\end{equation}
The \emph{Fisher spectral gap} is $\mu_1^{\mathrm{F}} > 0$, the smallest
positive eigenvalue of $-\Delta^{\mathrm{F}}$ on $L^2(\mathcal{S}, \gF)$.
\end{definition}

\begin{theorem}[Fisher spectral capacity sandwiching]
\label{thm:fisher-spectral-sandwich}
Let $(\mathcal{S}, \gF)$ be a compact statistical manifold of dimension $n$
with Ricci curvature $\mathrm{Ric}_{\gF} \ge (n-1)\kappa$ and diameter
$D_{\mathrm{F}} := \sup_{p,q} d_{\gF}(p,q)$.  The Fisher spectral gap
$\mu_1^{\mathrm{F}}$ satisfies:
\begin{enumerate}[label=\textnormal{(\roman*)}, leftmargin=*]
  \item \emph{(Cheeger lower bound.)}
  $\displaystyle \mu_1^{\mathrm{F}} \;\ge\; \frac{h_{\mathrm{F}}^2}{2}$,
  where $h_{\mathrm{F}}$ is the Cheeger isoperimetric constant of
  $(\mathcal{S}, \gF)$.
  \item \emph{(Lichnerowicz lower bound.)} If $\kappa > 0$, then
  $\mu_1^{\mathrm{F}} \ge n\kappa$.
  \item \emph{(Cheng upper bound.)}
  $\displaystyle \mu_1^{\mathrm{F}} \;\le\;
  \frac{C(n, \kappa)}{D_{\mathrm{F}}^2}$.
\end{enumerate}
These recover Branch~IX's Theorem~3.1 with $g = \gF$.  The Fisher spectral
capacity $C_{\mathrm{sp}}^{\mathrm{F}} := \mu_1^{\mathrm{F}} D_{\mathrm{F}}^2$
is therefore bounded by the same Cheeger--Ricci--Cheng sandwich as the parent
framework's spectral capacity.
\end{theorem}

\begin{proof}
These are standard results in spectral Riemannian geometry
(Cheeger 1970, Lichnerowicz 1958, Cheng 1975) applied to the specific metric
$\gF$.  No additional argument is needed beyond specialization.
\end{proof}

\subsection{Central conjecture}

\begin{mdframed}[backgroundcolor=gray!10, linewidth=0.4pt]
\begin{conjecture}[Fisher gap governs completion mixing]
\label{conj:fisher-mixing}
Let $(\mathcal{S}, \gF)$ be the Fisher-metric Davis manifold of a generative
model.  The spectral gap $\mu_1^{\mathrm{F}}$ of the Fisher Laplacian controls
the mixing time of the Markov chain on completions induced by autoregressive
sampling:
\begin{equation}\label{eq:mixing-bound}
t_{\mathrm{mix}}(\varepsilon)
\;\le\;
\frac{1}{\mu_1^{\mathrm{F}}}
\log\!\left(\frac{n}{\varepsilon}\right),
\end{equation}
where $n$ is the effective dimension of the completion space and $\varepsilon$
is the total-variation mixing tolerance.
\end{conjecture}
\end{mdframed}

This is the single central conjecture of Branch~X.  It is empirically testable:
compute $\mu_1^{\mathrm{F}}$ from the Fisher information of a transformer's
output distributions and check whether it predicts the number of sampling steps
needed for convergence to the stationary distribution over completions.

Motivation comes from the classical connection between spectral gaps and mixing
times on Riemannian manifolds (see e.g.\ Diaconis \& Stroock 1991): the
spectral gap of the generator of a reversible Markov chain controls the
exponential rate of convergence to equilibrium.  On the Fisher manifold, the
generator of the natural Brownian motion is precisely $\Delta^{\mathrm{F}}$,
and $\mu_1^{\mathrm{F}}$ controls its mixing.  The conjecture asserts that
autoregressive sampling, while not literally Brownian motion on
$(\mathcal{S}, \gF)$, is governed by the same spectral quantity
\emph{in the continuous-time diffusion limit} of the sampling chain
(or, equivalently, for any reversible approximation whose generator's
principal symbol agrees with $\Delta^{\mathrm{F}}$).

%====================================================
\section{Directions: dual connections and neural architecture}
\label{sec:directions}
%====================================================

This section collects conjectural connections between the $\alpha$-connection
family and the internal structure of transformer architectures.  These are
\emph{directions}, not results: they are clearly speculative, empirically
testable in several cases, and would require substantial additional work to
establish.  They are included to map the landscape of questions that
Branch~X's identification opens, not to claim answers.

\begin{conjecture}[Attention as $\alpha$-transport]
\label{conj:attention-transport}
Each attention head in a transformer performs approximate parallel transport of
information vectors along an effective $\alpha$-connection
$\nabla^{(\alpha_h)}$, where $\alpha_h$ depends on the head's learned
parameters:
\begin{itemize}
  \item Value projection transports content along the mixture connection
  $\nabla^{(m)}$, operating in expectation (mean) coordinates.
  \item Key-query interaction computes inner products in natural
  ($\theta$-)coordinates, where the exponential connection $\nabla^{(e)}$
  is flat.
  \item Multi-head combination approximates Levi-Civita transport
  $\nabla^{(0)}$ by averaging over heads with different effective $\alpha$.
\end{itemize}
Evidence required: empirical measurement of the effective $\alpha$ from
attention weight distributions across layers.
\end{conjecture}

\begin{conjecture}[Pythagorean decomposition]
\label{conj:pythagorean}
On an exponential family manifold, the generalized Pythagorean theorem for KL
divergence
\[
\DKL(p \| r) = \DKL(p \| q) + \DKL(q \| r)
\]
holds when $q$ is the $e$-projection of $p$ onto the $m$-geodesic through $r$
(Amari \& Nagaoka 2000, \S3.4).  Softmax normalization implements exactly this
$e$-projection, decomposing the total KL distance into orthogonal
information-geometric components.
\end{conjecture}

\begin{conjecture}[Layer norm as Fisher normalization]
\label{conj:layer-norm}
Layer normalization rescales activations to approximately unit Fisher norm,
acting as a gauge-fixing operation on the statistical manifold that prevents
the metric tensor from degenerating during deep composition.
\end{conjecture}

\begin{conjecture}[Residual stream as perturbed geodesic]
\label{conj:residual}
The residual connection $h_\ell = h_{\ell-1} + f_\ell(h_{\ell-1})$
implements a discretized geodesic equation on $(\mathcal{S}, \gF)$ with
layer-dependent perturbation.  The deviation from a true geodesic is
controlled by the Jacobi field equation on $(\mathcal{S}, \gF)$; reasoning
failures correspond to geodesic deviation exceeding the curvature-dependent
stability bound.
\end{conjecture}

Each of these conjectures is empirically testable but speculative at this
stage.  We record them as guide-posts for future work, not as claims of the
current paper.

%====================================================
\section{Cross-branch synthesis}
\label{sec:cross-branch}
%====================================================

The Fisher metric identification creates organic connections to each preceding
branch.  In each case, the connection arises from a classical identity in
information geometry, not from a forced analogy.

\subsection{Branch~VI: statistical mechanics}

\begin{theorem}[Fisher metric as free energy Hessian]
\label{thm:fisher-free-energy}
Let $Z(\theta) = \int h(x) \exp(\theta^i t_i(x))\, dx$ be the partition
function of an exponential family, and let
$\psi(\theta) = \log Z(\theta)$ be the log-partition function (free energy up
to temperature scaling).  Then:
\begin{equation}\label{eq:fisher-free-energy}
\gF_{ij}(\theta)
= \frac{\partial^2 \psi}{\partial\theta^i\, \partial\theta^j}
= \frac{\partial^2 (-F / T)}{\partial\theta^i\, \partial\theta^j},
\end{equation}
where $F = -T \psi(\theta)$ is the Helmholtz free energy at temperature $T$.
The Branch~VI susceptibility matrix
$\chi_{ij} = \partial^2 F / \partial h^i \partial h^j$ is (up to a
temperature factor) the Fisher metric at thermal equilibrium.
\end{theorem}

\begin{proof}
This is Proposition~\ref{prop:fisher-exp-family} restated in
thermodynamic language: $\gF = \nabla^2 \psi = \Cov[t_i, t_j]$ is both
the Fisher information and the fluctuation matrix of the sufficient
statistics, which in statistical mechanics is the susceptibility.
\end{proof}

\noindent\emph{Intuition.}  The Fisher metric at a point $\theta$ on the
statistical manifold measures how ``responsive'' the distribution $p_\theta$
is to small changes in $\theta$.  In thermodynamics, the susceptibility
$\chi$ measures how responsive a system's macroscopic properties are to
changes in external fields.  These are the same quantity because both ask:
``How much does the system fluctuate around its mean state?''  The
fluctuation--dissipation theorem of statistical mechanics is, in this light,
a statement about the Fisher metric: the equilibrium fluctuations (covariance
of sufficient statistics) equal the response to perturbation (Fisher
information).

\paragraph{Running example: HERALD at the antigenic transition.}
In HERALD's antigenic model, $\theta$ parameterizes the position in antigenic
space and $\psi(\theta)$ is the free energy of the lineage ensemble.  At
the Branch~VI completion transition ($\Gamma = 1$), the susceptibility
diverges: the system becomes maximally sensitive to perturbation.  In
Fisher-metric language, this means the Fisher information becomes large
(or ill-conditioned), and the Davis capacity bound
$C \le c / \sup K^{(0)}$ becomes tight.  The transition from ``single
dominant lineage'' to ``coexisting lineages'' is geometrically a transition
from a well-conditioned Fisher metric to an ill-conditioned one.

\begin{conjecture}[Phase transitions as Fisher degeneracy]
\label{conj:phase-transition}
The Branch~VI completion transition at $\Gamma = 1$ (the Davis trichotomy
boundary) corresponds to degeneracy of the Fisher metric:
$\det(\gF) \to 0$ as $\Gamma \to 1$.  The three regimes are:
\begin{itemize}
  \item Determined ($\Gamma > 1$): $\gF$ positive definite, unique completion
  geodesic, finite curvature.
  \item Critical ($\Gamma = 1$): $\gF$ degenerates along one or more
  directions, curvature diverges, phase transition.
  \item Underdetermined ($\Gamma < 1$): $\gF$ develops null directions,
  multiple completion geodesics, flat directions in the Fisher manifold.
\end{itemize}
\end{conjecture}

\subsection{Branch~VII: sheaf theory}

\begin{conjecture}[Obstruction as Fisher incompatibility]
\label{conj:sheaf-fisher}
The \v{C}ech obstruction class
$\delta(b) \in \check{H}^1(M, \mathcal{F}_{\mathrm{ext}})$ of Branch~VII
has an information-geometric interpretation: it measures the failure of local
Fisher metrics on constraint patches to glue into a global Fisher metric.
The obstruction vanishes if and only if local exponential families on each
patch of the constraint nerve are globally compatible---that is, there exists
a single exponential family on $M$ whose restriction to each patch agrees
with the local family up to controlled error.
\end{conjecture}

\subsection{Branch~VIII: renormalization group}

\begin{conjecture}[Data processing and the $C$-theorem]
\label{conj:data-processing}
The Branch~VIII $C$-theorem (completion entropy is monotonically
non-increasing under coarse-graining) is a consequence of the
\emph{data processing inequality for Fisher information}: for any Markov
kernel (coarse-graining map) $K$,
\[
\gF_{K \circ p}
\;\preceq\;
\gF_p
\]
in the Loewner order.  The Branch~VIII $C$-function is the trace (or
determinant) of the Fisher metric at scale $\ell$: it decreases under
each RG step because coarse-graining cannot increase Fisher information.
\end{conjecture}

\noindent\emph{Intuition.}  The data processing inequality says that
processing data through a noisy channel cannot increase Fisher information.
Coarse-graining in the RG sense is exactly such a channel: it maps
fine-grained distributions to coarse-grained ones.  Branch~VIII's
$C$-theorem says that a ``complexity'' measure decreases under RG flow.
The conjecture identifies this complexity measure as $\tr(\gF)$: the total
Fisher information at each scale.  This provides a statistical foundation
for the RG monotonicity that was postulated in Branch~VIII on physical
grounds.

\paragraph{Running example: AMC coarse-graining.}
In AMC's translator graph, coarse-graining corresponds to merging telemetry
formats that have ``similar enough'' translation behavior.  The data
processing inequality predicts that the Fisher information of the coarsened
graph is dominated by that of the fine-grained graph.  In concrete terms:
if AMC merges two telemetry formats whose reconstruction-error distributions
are close, the merged node's Fisher weight is less than or equal to the sum
of the original weights.  This controls the loss of discriminative power
under format consolidation.

\subsection{Branch~IX: spectral geometry}

\begin{theorem}[Fisher Laplacian on CSP graphs]
\label{thm:fisher-csp}
Let $G = (V, E)$ be a finite CSP graph with $n$ vertices and combinatorial
Laplacian $L_G$.  At each vertex $L_k$, the completion count $N_k := N_k(L_k)$
induces a scalar Fisher weight
\begin{equation}\label{eq:fisher-csp-weight}
\gF(L_k) = \frac{1}{N_k(L_k)},
\end{equation}
corresponding to the Fisher information of a uniform distribution over
$N_k$ compatible clue configurations.\footnote{If each solution $L_k$ admits
$N_k$ compatible clue configurations with equal probability $1/N_k$, the
Fisher information for the ``identity parameter'' is $1/N_k$.  More precisely,
this is the Fisher information for estimating which of $N_k$ equally likely
configurations generated the observed clues.
\emph{This is a choice of statistical model} (uniform likelihood over
completions), not a theorem.  The $1/N_k$ weighting is the simplest model
consistent with the completion-count structure; non-uniform priors over
clue configurations would yield different Fisher weights while preserving
the qualitative sandwiching structure of Theorem~\ref{thm:fisher-spectral-sandwich}.}

Define the diagonal Fisher weight matrix
$D_{\mathrm{F}} := \diag(\gF(L_1), \dots, \gF(L_n))$.  The
\emph{Fisher Laplacian on $G$} is
\begin{equation}\label{eq:fisher-laplacian-csp}
\Delta^{\mathrm{F}}_G
:= D_{\mathrm{F}}^{-1/2}\; L_G\; D_{\mathrm{F}}^{-1/2}.
\end{equation}
This is a symmetric positive semidefinite operator whose spectral gap
$\mu_1^{\mathrm{F}}(G)$ is the Fisher-weighted analog of the combinatorial
spectral gap $\mu_1(G)$ of Branch~IX.

The two spectral gaps are related by
\begin{equation}\label{eq:gap-relation}
\frac{\mu_1(G)}{\max_k \gF(L_k)}
\;\le\;
\mu_1^{\mathrm{F}}(G)
\;\le\;
\frac{\mu_1(G)}{\min_k \gF(L_k)},
\end{equation}
so the Fisher and combinatorial spectral gaps agree up to the condition number
of $D_{\mathrm{F}}$.
\end{theorem}

\begin{proof}
The matrix $\Delta^{\mathrm{F}}_G = D_{\mathrm{F}}^{-1/2} L_G
D_{\mathrm{F}}^{-1/2}$ is congruent to $L_G$ via $D_{\mathrm{F}}^{-1/2}$,
so it is symmetric and positive semidefinite with the same rank.  Its
eigenvalues are related to those of $L_G$ by the standard Rayleigh quotient
comparison: for any $f \perp \mathbf{1}$,
\[
\mu_1^{\mathrm{F}}
= \min_{f \perp \mathbf{1}} \frac{f^T \Delta^{\mathrm{F}}_G f}{f^T f}
= \min_{f \perp \mathbf{1}} \frac{(D_{\mathrm{F}}^{-1/2} f)^T L_G
(D_{\mathrm{F}}^{-1/2} f)}{f^T f}.
\]
Writing $h = D_{\mathrm{F}}^{-1/2} f$ and bounding the ratio
$\|f\|^2 / \|h\|^2$ by the extremal eigenvalues of $D_{\mathrm{F}}$
gives~\eqref{eq:gap-relation}.
\end{proof}

\begin{conjecture}[Fisher-weighted localization]
\label{conj:fisher-localization}
The Branch~IX result that completion count variation $N_k - \bar{N}$ lies in
the top eigenspace of $L_G$ has a Fisher-metric refinement: the
Fisher-weighted statistic
$\tilde{N}_k^{\mathrm{F}} := (N_k - \bar{N}) / \sqrt{\gF(L_k)}
= (N_k - \bar{N}) \sqrt{N_k}$
is an eigenvector (or near-eigenvector) of the Fisher Laplacian
$\Delta^{\mathrm{F}}_G$.
\end{conjecture}


%====================================================
\section{Finite-model computations}
\label{sec:finite-models}
%====================================================

We now compute the Fisher geometry on the two finite CSP graphs from
Branch~IX: the 288-vertex Shidoku graph and the 576-vertex Latin square graph.
These computations serve as sanity checks and as concrete instances where the
cross-branch predictions can be verified numerically.

\subsection{Fisher geometry of the Shidoku graph}

Recall from Branch~IX that the Shidoku graph $G_{\mathrm{Sh}}$ has
$n = 288$ vertices (the complete set of order-4 Shidoku solutions) and is
semiregular with two vertex classes determined by the Hamming-shell
completion count $N_k$ (the number of neighboring solutions at Hamming
distance exactly~$4$, corresponding to single-band swaps):
\begin{itemize}
  \item \textbf{Class A:} 192 vertices with $N_k = 4$.
  Fisher weight: $\gF = 1/4 = 0.25$.
  \item \textbf{Class B:} 96 vertices with $N_k = 8$.
  Fisher weight: $\gF = 1/8 = 0.125$.
\end{itemize}

The Fisher weight ratio is $\gF_A / \gF_B = 2$, so the condition number of
$D_{\mathrm{F}}$ is~$2$.  By~\eqref{eq:gap-relation}, the Fisher and
combinatorial spectral gaps differ by at most a factor of~$2$.

\paragraph{Worked computation: Fisher Laplacian eigenvalues.}
The combinatorial Laplacian $L_G$ of $G_{\mathrm{Sh}}$ has eigenvalues
$0 = \lambda_0 < \lambda_1 \le \cdots \le \lambda_{287}$.  The Fisher Laplacian
$\Delta^{\mathrm{F}} = D_{\mathrm{F}}^{-1/2}\, L_G\, D_{\mathrm{F}}^{-1/2}$
rescales each row and column by $(\gF(L_k))^{-1/2}$.  Since the graph has
only two Fisher-weight values, the rescaling acts blockwise:
\begin{itemize}[leftmargin=*]
  \item Edges between Class~A ($\gF = 1/4$) and Class~B ($\gF = 1/8$):
  rescaled by
  $(\gF_A \cdot \gF_B)^{-1/2} = (1/4 \cdot 1/8)^{-1/2} = \sqrt{32}
  \approx 5.66$.
  \item Edges within each class: rescaled by $(\gF_k)^{-1}$, i.e.\ by~$4$
  for Class~A and by~$8$ for Class~B.
  \item Diagonal entries: $d_i / \gF(L_i)$, preserving the Laplacian
  row-sum property.
\end{itemize}

Numerically, the combinatorial spectral gap is $\lambda_1 = 0.7085$ and the
Fisher spectral gap is $\mu_1^{\mathrm{F}} = 3.3394$.  The eigenvalue bounds
from Theorem~\ref{thm:fisher-csp} give:
\begin{align*}
\frac{\lambda_1}{\gF_{\max}}
&= \frac{0.7085}{0.25}
= 2.834
\;\le\; \mu_1^{\mathrm{F}} = 3.339, \\
\mu_1^{\mathrm{F}}
&\le\; \frac{\lambda_1}{\gF_{\min}}
= \frac{0.7085}{0.125}
= 5.668.
\end{align*}
The Fisher spectral gap is sandwiched between $2.834$ and $5.668$, a
factor-of-two window consistent with the condition number~$\kappa_{\mathrm{F}} = 2$.

\paragraph{Fisher-weighted completion count.}
The Fisher-weighted completion count statistic is
\[
\tilde{N}_k^{\mathrm{F}} := \frac{N_k - \bar{N}}{\sqrt{N_k}}
= \begin{cases}
(4 - 5.33)/\sqrt{4} = -1.33/2 \approx -0.67 & \text{(Class A, $N_k = 4$)}, \\
(8 - 5.33)/\sqrt{8} = 2.67/\sqrt{8} \approx 0.94 & \text{(Class B, $N_k = 8$)},
\end{cases}
\]
where $\bar{N} = (192 \times 4 + 96 \times 8)/288 = 5.33$.  The Fisher
weighting normalizes by $\sqrt{N_k}$, giving
more weight to vertices with fewer completions (higher Fisher information).
In raw units the signal amplification is
$(N_k - \bar{N}) \cdot \sqrt{N_k}$: Class~A emits $-2.67$ and Class~B
emits $+7.54$, so Class~B (fewer vertices, higher degree) is ``louder'' by
a factor of $7.54/2.67 \approx 2.8$ in Fisher-weighted coordinates.

\subsection{Fisher geometry of the Latin square graph}

The Latin square graph $G_{\mathrm{LS}}$ has $n = 576$ vertices (the complete
set of reduced $4 \times 4$ Latin squares, up to the appropriate equivalence)
and is again bipartite semiregular:
\begin{itemize}
  \item \textbf{Class A:} 432 vertices with $N_k = 72$.
  Fisher weight: $\gF = 1/72 \approx 0.0139$.
  \item \textbf{Class B:} 144 vertices with $N_k = 24$.
  Fisher weight: $\gF = 1/24 \approx 0.0417$.
\end{itemize}

The condition number of $D_{\mathrm{F}}$ is $\gF_B/\gF_A = (1/24)/(1/72) = 3$,
so the Fisher and combinatorial spectral gaps differ by at most a factor of~$3$.
The eigenvalue bounds give:
\begin{align*}
24\, \mu_1(G_{\mathrm{LS}})
&\;\le\; \mu_1^{\mathrm{F}}(G_{\mathrm{LS}})
\;\le\; 72\, \mu_1(G_{\mathrm{LS}}).
\end{align*}

\paragraph{Fisher-weighted intercalate count.}
The intercalate count $\mathcal{I}(L)$ is the variable that separates the two
classes in Branch~IX: $\mathcal{I} = 12$ for Class~A and $\mathcal{I} = 4$ for
Class~B.  The mean intercalate count is
$\bar{\mathcal{I}} = (432 \times 12 + 144 \times 4)/576 = 10$.  The
Fisher-weighted intercalate statistic is
\[
\tilde{\mathcal{I}}^{\mathrm{F}}(L) :=
\frac{\mathcal{I}(L) - \bar{\mathcal{I}}}{\sqrt{\gF(L)}}
= (\mathcal{I}(L) - \bar{\mathcal{I}}) \cdot \sqrt{N_k(L)}
= \begin{cases}
(12 - 10)\sqrt{72} = 2\sqrt{72} \approx 16.97 & \text{(Class A)}, \\
(4 - 10)\sqrt{24} = -6\sqrt{24} \approx -29.39 & \text{(Class B)}.
\end{cases}
\]
The Fisher weighting \emph{amplifies} the separation for Class~B vertices
(which have fewer completions and hence higher Fisher information),
producing an asymmetric signal: Class~B is ``louder'' in Fisher-weighted
coordinates because it has higher information density per vertex.  The
absolute separation is
$|\tilde{\mathcal{I}}^{\mathrm{F}}_A| + |\tilde{\mathcal{I}}^{\mathrm{F}}_B|
\approx 46.4$,
compared to the unweighted separation $|2| + |{-}6| = 8$.  Fisher weighting
amplifies the signal by a factor of approximately~$5.8$.

\paragraph{Comparison of Fisher weighting across the two graphs.}
\begin{table}[htbp]
\centering
\caption{\textbf{Fisher geometry comparison} of the Shidoku and Latin square
CSP graphs.  Condition number $\kappa_{\mathrm{F}}$ controls the gap between
Fisher and combinatorial spectral gaps.}
\label{tab:fisher-comparison}
\begin{tabular}{lccccc}
\toprule
Graph & $n$ & $\kappa_{\mathrm{F}}$ & $\gF_{\min}$ & $\gF_{\max}$ &
Fisher signal amplification \\
\midrule
Shidoku & 288 & 2 & $1/8$ & $1/4$ & $\times 2.8$ \\
Latin square & 576 & 3 & $1/72$ & $1/24$ & $\times 5.8$ \\
\bottomrule
\end{tabular}
\end{table}

The Latin square graph has both a higher condition number and stronger Fisher
signal amplification, reflecting its more heterogeneous completion structure.
This predicts that Fisher-weighted diagnostics will be more informative on the
Latin square graph than on the Shidoku graph---exactly the opposite of what
one might expect from the unweighted combinatorial analysis alone.

\subsection{Partition function cross-check with Branch~VI}

As a cross-check with the statistical mechanics formulation of Branch~VI, we
verify the identity $\gF(\beta) = \partial^2 \psi / \partial\beta^2$ on the
Shidoku model.

\paragraph{Setup.}  Branch~VI defines a partition function over Shidoku clue
configurations at inverse temperature $\beta$:
\[
Z(\beta) = \sum_{\gamma \in \Gamma} \exp(-\beta\, E[\gamma]),
\]
where $\Gamma$ is the set of clue configurations and $E[\gamma]$ is an energy
functional measuring the constraint-violation count.  The log-partition function
is $\psi(\beta) = \log Z(\beta)$, and the Branch~VI specific heat is
\[
C_V(\beta) = \beta^2\,\bigl(\langle E^2 \rangle_\beta -
\langle E \rangle_\beta^2\bigr)
= \beta^2\, \Var_\beta[E].
\]

\paragraph{Information-geometric identity.}
Treating $\beta$ as a one-dimensional natural parameter with sufficient
statistic $t(\gamma) = -E[\gamma]$, the Boltzmann distribution
$p_\beta(\gamma) = \exp(-\beta E[\gamma])/Z(\beta)$
is a one-dimensional exponential family.  By
Proposition~\ref{prop:fisher-exp-family}, the Fisher metric at $\beta$ is
\[
\gF(\beta) = \frac{\partial^2 \psi}{\partial\beta^2}
= \Var_\beta[E] = \frac{C_V(\beta)}{\beta^2}.
\]

\paragraph{Verification protocol.}
\begin{enumerate}[label=\textnormal{(\arabic*)}, leftmargin=*]
  \item Compute $Z(\beta)$ by exact enumeration over Shidoku clue
  configurations (feasible because $|\Gamma| = 288 \times
  \binom{16}{k}$ for each clue count $k$, which is computationally accessible
  for $k \le 6$).
  \item Compute $C_V(\beta)$ from the thermal ensemble by finite differencing
  of $\psi(\beta)$ or by direct variance computation.
  \item Compute $\gF(\beta)$ from the information-geometric definition as
  $\partial^2 \psi / \partial\beta^2$.
  \item Verify numerical agreement: $C_V(\beta)/\beta^2 = \gF(\beta)$ to
  machine precision.
\end{enumerate}

This computation is a \emph{non-trivial cross-branch consistency check}: it
verifies that the same function $\psi = \log Z$ governs both the Branch~VI
thermodynamics and the Branch~X information geometry.  Agreement confirms that
the Fisher metric identification is consistent with the statistical-mechanics
framework.  Disagreement would indicate a bug in the formalism.

\paragraph{Expected behavior near the completion transition.}
At the Branch~VI completion transition ($\Gamma = 1$, the Davis trichotomy
boundary), the specific heat $C_V(\beta)$ diverges (or peaks sharply on a
finite system), corresponding to a large Fisher metric $\gF(\beta)$.  This is
exactly the prediction of Conjecture~\ref{conj:phase-transition}: the Fisher
metric becomes large near the transition, and its condition number
deteriorates, signaling the onset of multi-valued completions.


\subsection{Computational verification}
\label{sec:computational-evidence}

The claims of Sections~\ref{sec:statistical-manifolds}--\ref{sec:finite-models} and the
cross-branch identity of Theorem~\ref{thm:fisher-free-energy} were verified
numerically by an independent Python script that takes no analytic shortcuts:
every eigenvalue, variance, and gradient is computed from scratch on the
288-vertex Shidoku model.  We report the results below; the full script and
raw output are available in the supplementary materials.\footnote{Script:
\texttt{branch\_x\_evidence.py}.  All tests pass on NumPy~1.26 /
SciPy~1.14 with no manual tolerances wider than $10^{-4}$.}

\paragraph{Part~1: Softmax Fisher geometry (Proposition~\ref{prop:fisher-exp-family} through Theorem~\ref{thm:curvature-capacity}).}
For $\bm{\theta} = (2, 1, 0.5, 0)$ with $V = 4$:
\begin{enumerate}[label=\textnormal{(\roman*)}, leftmargin=*]
  \item $\gF = \diag(p) - pp^\top$ confirmed symmetric, PSD, rank $V{-}1 = 3$,
  with null vector $\mathbf{1}$ (shift invariance).
  Eigenvalues: $0.3464, 0.1590, 0.0908, 0.0000$.  Condition number
  $\kappa = 3.81 < 10$.
  \item Numerical Hessian of $\psi(\bm{\theta}) = \log \sum_v e^{\theta_v}$
  agrees with $\gF$ to $1.93 \times 10^{-6}$ (central differences,
  $\varepsilon = 10^{-5}$).
  \item Legendre duality:
  $\bm{\eta} = \nabla\psi = \mathrm{softmax}(\bm{\theta})$ verified to
  $10^{-8}$; conjugate potential $\varphi(\bm{\eta}) = \bm{\theta} \cdot
  \bm{\eta} - \psi$ matches negative entropy to $10^{-10}$.
  \item Fisher--Rao distance $d_{\mathrm{FR}}(p, q)$ versus
  $\sqrt{2\,\DKL(p \| q)}$: ratio converges to $1.0000$ at
  $\varepsilon = 10^{-3}$ and remains above $0.97$ at $\varepsilon = 0.5$.
  \item Cram\'{e}r--Rao saturation: $50{,}000$-sample MLE over $1{,}000$
  trials gives $n \cdot \Cov[\hat{p}] / \gF$ with mean diagonal ratio
  $0.983$, confirming that the categorical MLE saturates the bound.
\end{enumerate}

\begin{figure}[htbp]
\centering
\begin{minipage}[t]{0.38\textwidth}
\centering
\includegraphics[width=\textwidth]{fig_01_fisher_heatmap.png}
\end{minipage}
\hfill
\begin{minipage}[t]{0.58\textwidth}
\centering
\includegraphics[width=\textwidth]{fig_02_fisher_rao_kl.png}
\end{minipage}
\caption{\textbf{Left:} Fisher information metric $\gF = \diag(p) - pp^\top$
for $p = \mathrm{softmax}(2, 1, 0.5, 0)$.  Diagonal entries decrease with
probability; off-diagonal entries are negative (anti-correlated).
\textbf{Right:} Fisher--Rao distance versus $\sqrt{2\,\DKL}$ as a function
of perturbation scale $\varepsilon$.  The two metrics agree to second order
(ratio $\to 1$ as $\varepsilon \to 0$), confirming
Theorem~\ref{thm:kl-fisher}.}
\label{fig:fisher-geometry}
\end{figure}

\paragraph{Part~2: Shidoku Fisher Laplacian (Theorem~\ref{thm:fisher-spectral-sandwich}).}
\begin{enumerate}[label=\textnormal{(\roman*)}, leftmargin=*]
  \item Backtracking enumeration finds exactly $288$ Shidoku solutions.
  \item Adjacency at Hamming distance exactly~$4$ (a single-band swap) yields
  $768$ edges on $1$ connected component, with degree range $[4, 8]$.
  \item Hamming-shell completion counts $N_k$ (the number of solutions at
  distance exactly~$4$ from each vertex) produce two Fisher-weight classes:
  \begin{center}
  \begin{tabular}{cccl}
  \toprule
  $N_k$ & vertices & $\gF = 1/N_k$ & class \\
  \midrule
  $4$ & $192$ & $0.2500$ & high-information \\
  $8$ & $96$  & $0.1250$ & low-information \\
  \bottomrule
  \end{tabular}
  \end{center}
  Condition number $\kappa_{\mathrm{F}} = 2.00$.
  \item Combinatorial spectral gap $\lambda_1 = 0.7085$.
  Fisher Laplacian spectral gap $\mu_1^{\mathrm{F}} = 3.3394$.
  Sandwiching inequality:
  \[
  \underbrace{\frac{\lambda_1}{\gF_{\max}}}_{2.8340}
  \;\le\; \mu_1^{\mathrm{F}} = 3.3394
  \;\le\; \underbrace{\frac{\lambda_1}{\gF_{\min}}}_{5.6680}.
  \]
  \textbf{Theorem~\ref{thm:fisher-spectral-sandwich} confirmed.}
\end{enumerate}

\begin{figure}[htbp]
\centering
\includegraphics[width=0.75\textwidth]{fig_04_sandwich.png}
\caption{\textbf{Spectral gap sandwiching (Theorem~\ref{thm:fisher-spectral-sandwich}).}
The combinatorial spectral gap $\lambda_1 = 0.708$ (dotted white) is
amplified by the Fisher weighting to $\mu_1^{\mathrm{F}} = 3.339$ (solid
gold), landing within the sandwich bounds $[2.834, 5.668]$ (blue band).
The $\kappa_{\mathrm{F}} = 2$ condition number gives a tight factor-of-two
window.}
\label{fig:spectral-sandwich}
\end{figure}

\paragraph{Part~3: Partition function cross-check (Theorem~\ref{thm:fisher-free-energy}).}
Using energy $E = -\sum_{r} s_{rr}$ (negative diagonal sum), the $288$
solutions partition into a symmetric energy spectrum with nine distinct levels
($E \in \{-14, \ldots, -6\}$).  At each inverse temperature
$\beta \in [0.05, 4.0]$:
\begin{enumerate}[label=\textnormal{(\roman*)}, leftmargin=*]
  \item $\Var_\beta[E]$ from the Boltzmann ensemble agrees with
  $\partial^2\psi/\partial\beta^2$ (numerical second derivative) to four
  significant figures across the full range.  Maximum absolute deviation:
  $1.39 \times 10^{-4}$.
  Representative values:
  \begin{center}
  \begin{tabular}{cccc}
  \toprule
  $\beta$ & $\Var_\beta[E]$ & $\psi''(\beta)$ & ratio \\
  \midrule
  $0.10$ & $3.3071$ & $3.3069$ & $0.9999$ \\
  $1.00$ & $1.6880$ & $1.6881$ & $1.0000$ \\
  $3.00$ & $0.1278$ & $0.1278$ & $1.0001$ \\
  \bottomrule
  \end{tabular}
  \end{center}
  \item The fluctuation-dissipation theorem $C_V = \beta^2 \Var_\beta[E]$
  holds to machine precision ($< 10^{-15}$), and $C_V$ computed from
  $-\partial\langle E\rangle / \partial\beta$ agrees with
  $\beta^2 \Var_\beta[E]$ to $10^{-4}$.
  \item The Fisher metric peaks at $\beta \approx 0.05$ (maximum fluctuations),
  while $C_V$ peaks at $\beta \approx 1.51$, consistent with the predicted
  phase-transition signature.
\end{enumerate}

\begin{figure}[htbp]
\centering
\includegraphics[width=0.85\textwidth]{fig_03_partition.png}
\caption{\textbf{Left:} Fisher metric $\gF(\beta) = \Var_\beta[E]$ (blue)
overlaid with $\psi''(\beta)$ (pink dashed) on the 288-solution Shidoku
ensemble.  The two curves are indistinguishable, confirming
$\gF = \partial^2\psi/\partial\beta^2$
(Theorem~\ref{thm:fisher-free-energy}).
\textbf{Right:} Specific heat $C_V(\beta) = \beta^2 \Var_\beta[E]$, peaking
at $\beta \approx 1.50$ (the completion transition).  The Fisher metric
peaks at maximum fluctuations ($\beta \to 0$), consistent with
Conjecture~\ref{conj:phase-transition}.}
\label{fig:partition-function}
\end{figure}

\paragraph{Part~4: Supplementary verifications.}
\begin{enumerate}[label=\textnormal{(\roman*)}, leftmargin=*]
  \item \emph{von Mises--Fisher (VIDAR).}  For $d = 128$, the ratio
  $A_d(\kappa) = I_{d/2}(\kappa) / I_{d/2-1}(\kappa)$ approaches~$1$ as
  $\kappa$ increases ($A_d(1000) = 0.938$), confirming
  $\gF \to \kappa \cdot g_{\mathrm{round}}$ on $S^{d-1}$.
  \item \emph{Data processing inequality.}  Progressive coarsening
  (merging categories) on a $V{=}5$ softmax distribution yields strictly
  decreasing $\tr(\gF)$: $0.565 \to 0.562 \to 0.549 \to 0.475$.
  \item \emph{Amari--Chentsov tensor.}  $T_{ijk}$ computed analytically for
  $V{=}4$ is totally symmetric with $\norm{T}_F = 0.254 > 0$, confirming
  nonzero Levi-Civita curvature on non-uniform categorical distributions
  (while the $e$- and $m$-connections remain individually flat, per
  Proposition~\ref{def:alpha-connection}).
\end{enumerate}

\begin{figure}[htbp]
\centering
\begin{minipage}[t]{0.48\textwidth}
\centering
\includegraphics[width=\textwidth]{fig_05_vmf.png}
\end{minipage}
\hfill
\begin{minipage}[t]{0.48\textwidth}
\centering
\includegraphics[width=\textwidth]{fig_06_dpi.png}
\end{minipage}
\caption{\textbf{Left:} Convergence of the von Mises--Fisher ratio
$A_d(\kappa) \to 1$ across embedding dimensions $d = 16, 64, 128, 512$.
Higher dimensions require larger concentration $\kappa$ to approach the
round-metric limit, but all curves converge, confirming
$\gF \to \kappa \cdot g_{\mathrm{round}}$ on $S^{d-1}$.
\textbf{Right:} Data processing inequality under progressive coarsening
($V = 5 \to 2$).  $\tr(\gF)$ decreases monotonically (left panel);
the pseudo-determinant $\det^*(\gF)$ increases (right panel, log scale)
because the surviving eigenvalues become less degenerate as categories merge.}
\label{fig:vmf-dpi}
\end{figure}

All eleven tested claims pass.  The script provides reproducible
computational evidence that the information-geometric framework of
Branch~X is internally consistent and that the cross-branch identities
with Branch~VI (partition functions) and Branch~IX (spectral geometry)
hold numerically.

%====================================================
\section{The empirical program}
\label{sec:empirical}
%====================================================

The Fisher metric is computable from data.  This section describes how to
measure it on real transformers, what predictions Branch~X makes, and what
success or failure would mean.

\subsection{Measurement protocol}

The following algorithm describes how to extract Fisher-geometric observables
from a transformer at inference time.  It is designed to parallel the distortion
audit protocol (Algorithm~1 in \cite{davis2025manifold}): where that protocol
measures geodesic-Euclidean distortion ratios $\rho_j$ along paths, this
protocol measures Fisher-Rao distances, curvature, and spectral gap along
completion paths.

\begin{algorithm}[H]
\caption{Fisher-Metric Extraction from Transformer Logits}
\label{alg:fisher-extraction}
\begin{algorithmic}[1]
\Require Transformer with $L$ layers, vocabulary $\mathcal{V}$ ($|\mathcal{V}|=V$),
completion tasks $\{(\text{prompt}_k, \text{completion}_k)\}_{k=1}^{N_{\mathrm{task}}}$
\Ensure Fisher-Rao distances $\{d^{\mathrm{F}}_k\}$, curvature profiles
$\{K_k(\ell)\}$, spectral gap estimate $\hat\mu_1^{\mathrm{F}}$
\For{each task $k = 1, \dots, N_{\mathrm{task}}$}
  \State Feed $\text{prompt}_k$ through the transformer
  \For{each layer $\ell = 1, \dots, L$}
    \State Extract logit vector $z_{\ell} \in \RR^V$
    \State Compute softmax: $p_{\ell}^v = \exp(z_\ell^v)/\sum_w \exp(z_\ell^w)$
    \State \textbf{Fisher metric:}
    $\gF(\ell) = \diag(p_\ell) - p_\ell\, p_\ell^\top$
    \Comment{Eq.~\eqref{eq:softmax-fisher-explicit}}
    \State Store eigenvalues $\lambda_1(\ell) \ge \cdots \ge \lambda_{V-1}(\ell) > 0$
    of $\gF(\ell)$
  \EndFor
  \State \textbf{Fisher-Rao path length:}
  $d_k^{\mathrm{F}} = \sum_{\ell=1}^{L-1}
  \sqrt{(z_{\ell+1} - z_\ell)^\top\, \gF(\ell)\, (z_{\ell+1} - z_\ell)}$
  \State \textbf{Layer-wise curvature proxy:}
  $K_k(\ell) = \kappa(\gF(\ell)) := \lambda_1(\ell)/\lambda_{V-1}(\ell)$
  \Comment{Condition number}
  \State \textbf{KL divergence:}
  $D_k(\ell) = \DKL(p_\ell \| p_{\ell+1})$
  for each $\ell$
\EndFor
\State \textbf{Aggregate Fisher-Rao statistics:}
mean $\bar{d}^{\mathrm{F}}$, std $\sigma_d$, percentiles
\State \textbf{Spectral gap estimate:}
$\hat\mu_1^{\mathrm{F}} = \text{median over tasks of } \min_\ell \lambda_{V-1}(\ell)$
\State \Return $\{d_k^{\mathrm{F}}\}$, $\{K_k(\ell)\}$, $\hat\mu_1^{\mathrm{F}}$,
$\{D_k(\ell)\}$
\end{algorithmic}
\end{algorithm}

\paragraph{Implementation notes.}
\begin{itemize}[leftmargin=*]
  \item \textbf{Computational cost.}  The Fisher metric
  $\gF = \diag(p) - pp^\top$ is computed from the softmax output without
  additional forward passes.  Its eigendecomposition costs $O(V^2)$ per layer
  per position, which is dominated by the transformer's $O(V \cdot d_{\mathrm{model}})$
  output projection.  For large vocabularies ($V \sim 50{,}000$), the full
  eigendecomposition can be replaced by a rank-$k$ approximation using
  randomized SVD, since the Fisher matrix of a peaked softmax has rapidly
  decaying eigenvalues.

  \item \textbf{Sherman--Morrison structure.}
  The matrix $\gF = \diag(p) - pp^\top$ is a rank-1 perturbation of a
  diagonal matrix.  The inverse is available in closed form via the
  Sherman--Morrison formula:
  \[
  (\gF)^{-1} = \diag(p)^{-1} + \frac{\diag(p)^{-1}\, pp^\top\, \diag(p)^{-1}}
  {1 - p^\top \diag(p)^{-1} p}
  = \diag(1/p_i) + \mathbf{1}\mathbf{1}^\top,
  \]
  where the second equality holds on the $(V{-}1)$-dimensional subspace
  orthogonal to $\mathbf{1}$.  This allows $O(V)$ computation of quadratic
  forms involving $(\gF)^{-1}$.

  \item \textbf{Batch processing.}  For a batch of $B$ prompts, Fisher metrics
  at all layers can be computed in parallel, since each depends only on the
  local softmax output.

  \item \textbf{Low-rank Fisher for hidden states.}
  If Fisher information is desired for internal representations (not just the
  output softmax), the empirical Fisher
  $\hat{\gF} = \frac{1}{N}\sum_{n=1}^N
  \nabla_\theta \ell_n \,(\nabla_\theta \ell_n)^\top$
  can be estimated via gradient samples; low-rank approximations such as
  K-FAC~\cite{martens2020} reduce this to $O(d_{\mathrm{model}}^2)$ per layer.
\end{itemize}

\begin{algorithm}[H]
\caption{Fisher Curvature Diagnostic for Completion Tasks}
\label{alg:fisher-diagnostic}
\begin{algorithmic}[1]
\Require Fisher-Rao distances $\{d_k^{\mathrm{F}}\}$, curvature profiles
$\{K_k(\ell)\}$, KL divergences $\{D_k(\ell)\}$, task difficulty labels
$\{y_k\}$ (correct/incorrect completion)
\Ensure Correlation analysis, curvature alarm thresholds
\State \textbf{Hypothesis 1 (Fisher-difficulty correlation):}
\State \quad Compute rank correlation $\rho(d^{\mathrm{F}}, y)$; report $p$-value
\State \textbf{Hypothesis 2 (curvature predicts failure):}
\State \quad For each task $k$, compute
$K_k^{\max} := \max_\ell K_k(\ell)$
\State \quad Fit logistic regression: $\Prob(y_k = \text{incorrect})
\sim \beta_0 + \beta_1 K_k^{\max}$
\State \quad Report coefficient $\beta_1$ and significance
\State \textbf{Hypothesis 3 (KL budget predicts failure):}
\State \quad Compute total KL: $\tau_k := \sum_\ell D_k(\ell)$
\State \quad Test whether failures cluster among tasks with $\tau_k > \tau_{\mathrm{thresh}}$
\State \textbf{Curvature alarm:}
\State \quad Set $K_{\mathrm{alarm}} = $ 90th percentile of $K_k^{\max}$
among correct completions
\State \quad Flag tasks with $K_k^{\max} > K_{\mathrm{alarm}}$ for abstention
\State \Return Correlations, thresholds, alarm rate
\end{algorithmic}
\end{algorithm}

\paragraph{Decision rules.}
Table~\ref{tab:fisher-diagnostic-rules} summarizes the diagnostic decision rules
analogous to the distortion audit gates in \cite{davis2025manifold},
Table~6.

\begin{table}[htbp]
\centering
\caption{\textbf{Fisher curvature diagnostic decision rules.}  Thresholds are
illustrative; actual values depend on architecture and task.}
\label{tab:fisher-diagnostic-rules}
\begin{tabularx}{\linewidth}{@{}p{0.28\linewidth}p{0.30\linewidth}p{0.34\linewidth}@{}}
\toprule
\textbf{Condition} & \textbf{Interpretation} & \textbf{Action} \\
\midrule
$K_k^{\max} < K_{\mathrm{nominal}}$ &
Low curvature; Fisher geometry reliable &
Trust completion; report $d^{\mathrm{F}}_k$ as confidence proxy \\
\addlinespace[2pt]
$K_{\mathrm{nominal}} \le K_k^{\max} < K_{\mathrm{alarm}}$ &
Moderate curvature; capacity may be reduced &
Report with reduced confidence; increase
sampling temperature or beam width \\
\addlinespace[2pt]
$K_k^{\max} \ge K_{\mathrm{alarm}}$ &
High curvature; capacity bound may be violated &
\textbf{Abstain} or defer to ensemble; flag for human review \\
\addlinespace[2pt]
$\tau_k > \tau_{\mathrm{max}}$ &
KL budget exceeded; accumulated information loss too large &
Completion path has exceeded tolerance; truncate or restart \\
\bottomrule
\end{tabularx}
\end{table}

\subsection{What success would mean}

If the empirical program confirms the predictions of Branch~X:
\begin{itemize}
  \item The Davis framework's curvature bounds become \emph{measurable
  diagnostics} for transformer reliability, computable at inference time.
  \item Reasoning failures become \emph{predictable} from Fisher-geometric
  invariants before they occur: a ``curvature alarm'' that fires when the
  accumulated Fisher-Rao distance or sectional curvature exceeds
  tolerance.
  \item The tolerance budget $\tau$ becomes a \emph{computable quantity},
  not a fitted parameter: it is determined by the Fisher geometry of the
  output distribution and the acceptable KL divergence.
\end{itemize}

\paragraph{Running example: HERALD success scenario.}
If Fisher curvature along epitope-restricted mutation paths correlates with
antigenic escape risk, then the HERALD system could use
Algorithm~\ref{alg:fisher-diagnostic} as a \emph{pre-screening} tool: compute
the Fisher-Rao distance between a new variant and the vaccine anchor in the
lineage posterior space, estimate the accumulated curvature, and flag variants
whose Fisher-geometric trajectory enters the high-curvature regime as escape
candidates \emph{before} neutralization assay data are available.  This would
provide earlier warning than the current protocol (which requires wet-lab assay
results to update the drift statistic $D(t)$).

\paragraph{Running example: VIDAR success scenario.}
If Fisher curvature of the softmax output distributions during identity
embedding predicts deepfake detection difficulty, then VIDAR could add a
Fisher-based pre-filter: compute $K_k^{\max}$ for each clip, and for clips
where $K_k^{\max} > K_{\mathrm{alarm}}$, escalate to ensemble methods or
human review before running the full RIM feature pipeline.  The curvature
alarm would act as an uncertainty quantification layer that is \emph{internal}
to the model's geometry, rather than bolted on post hoc.

\subsection{What failure would mean}

If the empirical program fails:
\begin{itemize}
  \item The Fisher metric may not be the right instantiation of the Davis
  metric: perhaps the transformer's internal geometry is not
  well-approximated by its \emph{output} statistics.  The Fisher metric
  of intermediate representations (hidden states, attention patterns) might
  be more appropriate than the Fisher metric of the final softmax.
  \item The exponential family assumption may be too restrictive: sparse
  attention, mixture-of-experts, and other non-softmax output heads
  produce non-exponential distributions for which the dual-flatness results
  of Section~\ref{sec:alpha-connections} do not hold.
  \item The finite-dimensional parameterization may be inadequate: for
  transformers with billions of parameters, the full Fisher matrix is
  intractable, and low-rank approximations (K-FAC, diagonal Fisher,
  empirical Fisher) introduce systematic error that could invalidate
  the curvature predictions.
\end{itemize}
Each failure mode is informative, falsifiable, and points toward a specific
refinement.  This is a feature of the framework, not a limitation: it is
designed to fail in scientifically useful ways.

\paragraph{Failure mode analysis.}
Table~\ref{tab:failure-modes} maps each failure mode to its diagnostic
signature and the corrective action it would motivate.

\begin{table}[htbp]
\centering
\caption{\textbf{Failure mode analysis for the empirical program.}}
\label{tab:failure-modes}
\begin{tabularx}{\linewidth}{@{}p{0.25\linewidth}p{0.32\linewidth}p{0.35\linewidth}@{}}
\toprule
\textbf{Failure mode} & \textbf{Diagnostic signature} & \textbf{Corrective action} \\
\midrule
Output Fisher insufficient &
No correlation between $d^{\mathrm{F}}$ and task difficulty; curvature alarm
has no predictive power &
Compute Fisher metric on hidden states via empirical Fisher or K-FAC; test
intermediate-layer Fisher \\
\addlinespace[2pt]
Non-exponential outputs &
Dual flatness violated; Cram\'er--Rao not saturated; capacity bounds loose &
Use general $\alpha$-geometry on curved manifolds; relax exponential family
assumption IG-A3 \\
\addlinespace[2pt]
Low-rank approximation error &
Fisher spectrum poorly estimated; eigenvalue ratios unstable across runs &
Increase rank in K-FAC/randomized SVD; use full Fisher on smaller models
as ground truth \\
\addlinespace[2pt]
Curvature--failure decorrelation &
Failures occur at low curvature; successes occur at high curvature &
Fisher metric captures wrong notion of difficulty; consider non-geometric
failure modes (data quality, distribution shift) \\
\bottomrule
\end{tabularx}
\end{table}

%====================================================
\section{Limitations and open problems}
\label{sec:limitations}
%====================================================

\subsection{Technical limitations}

\paragraph{Scope of the metric identification.}
The identification of the Fisher metric as the unique Davis metric
(Corollary~\ref{cor:canonical-davis}) is conditional on axioms IG-A1--A4:
that the Davis manifold carries statistical-model structure and that the
metric is required to be invariant under sufficient-statistic reductions.
Outside this statistical-manifold regime---for example, on Davis manifolds
arising from purely geometric or topological sameness structures that do not
naturally parameterize probability distributions---metric selection remains
an open problem.  The results of this branch apply whenever the axioms hold,
but do not claim to resolve metric selection in full generality.

\paragraph{Exponential family scope.}
The cleanest results of this branch---dual flatness
(Theorem~\ref{thm:e-flat}), Cram\'er--Rao saturation
(Corollary~\ref{cor:saturation}), the Pythagorean theorem
(Conjecture~\ref{conj:pythagorean})---hold only on exponential families.
Non-exponential output distributions require $\alpha$-geometry on general
curved statistical manifolds, where computations are harder and uniqueness
results (beyond \v{C}encov's, which still holds) are weaker.

\paragraph{Discrete-to-continuous gap.}
The Shidoku and Latin square computations of Section~\ref{sec:finite-models}
are on discrete CSP graphs with finitely many vertices.  The continuous
statistical manifold theory of Sections~\ref{sec:statistical-manifolds}--\ref{sec:curvature-capacity} applies in the limit of large, smoothly
parameterized families.  Bridging the gap requires scaling analysis analogous
to Branch~IX's generalization boundary (the point at which finite-model
spectral results extend to families of CSP instances).

\paragraph{$\alpha$-observability.}
Conjectures~\ref{conj:attention-transport}--\ref{conj:residual} posit specific
$\alpha$ values for transformer components (attention heads, layer norm,
residual stream).  These are empirically testable but theoretically
unmotivated at present: we have no architectural argument for why a particular
component should implement transport along a particular $\alpha$-connection.
The effective $\alpha$ may vary across heads, layers, and inputs.

\subsection{Relation to existing literature}

Information geometry of neural networks is an active and mature field.
Amari's natural gradient~\cite{amari1998} uses the Fisher metric to define
steepest descent in distribution space.
Martens~\cite{martens2020} provides a comprehensive analysis of the natural
gradient and its approximations.
Kunstner et al.~\cite{kunstner2019} study the gap between the exact and
empirical Fisher, which is directly relevant to the tractability of
Branch~X's measurement program.

This branch's contribution is not the Fisher metric per se, but:
\begin{itemize}
  \item its identification as the unique Davis metric (within the
  statistical-manifold regime) via \v{C}encov;
  \item the dual-connection extension of the parent framework;
  \item the cross-branch synthesis (Fisher = free energy Hessian =
  sheaf compatibility = RG monotone = spectral weighting);
  \item the specific empirical measurement program connecting abstract
  geometry to transformer observables.
\end{itemize}
The branch should be understood as building on, not competing with, the
existing information geometry literature.

\subsection{Connection to the Geometry of Sameness}

In the Geometry of Sameness~\cite{davis2025sameness}, the manifold-based
semantic realization (MSR, Definition~3) requires a Riemannian metric $g$
on the semantic manifold $\mathcal{M}$ but does not specify how to choose it.
Corollary~\ref{cor:canonical-davis} resolves this: when the feature spaces
$\{V_i\}$ of the MSR parameterize probability distributions, the Fisher
metric is the canonical choice.

\begin{remark}[\v{C}encov and the metric-selection problem]
\label{rmk:cencov-sameness}
The Geometry of Sameness defines $\varepsilon$-equivalence of realization
categories $\SamTrans(S)$ and $\SamGeom(S)$ and constructs functors $F_S$,
$G_S$ between them.  The functor $F_S$ builds a manifold from a translator
graph but does not specify which metric to place on the resulting manifold
beyond requiring bounded distortion.

\v{C}encov's theorem provides the missing prescription.  When the translator
graph has a statistical interpretation---each node parameterizes a
distribution over reconstruction errors across formats---the Fisher metric is
the unique metric on the resulting manifold that is invariant under sufficient
statistic reductions.  This means $F_S$ should use the Fisher metric, and
$G_S$ should use the inverse Fisher metric to build transition translators.

A full functor $\Phi_S : \SamGeom^0(S) \to \SamInfo^0(S)$ extending the
$\varepsilon$-equivalence chain to include information-geometric realizations
is a natural next step.  Such a functor would compose with $F_S$ and $G_S$
to form a triangle of $\varepsilon$-equivalences:
\[
\SamTrans^0(S) \xrightarrow{F_S} \SamGeom^0(S) \xrightarrow{\Phi_S}
\SamInfo^0(S) \xrightarrow{G'_S} \SamTrans^0(S).
\]
We defer the construction of $\Phi_S$ to a future synthesis paper; the
present branch establishes the identification $g = c \cdot \gF$ that would
form its core content.
\end{remark}

The Geometry of Sameness's smooth chartability condition (Definition~8) is
implied by \emph{Fisher regularity}: the condition that
$\gF(\theta)$ is positive definite with bounded condition number on the
in-regime set $\Omega^{\mathrm{ideal}}$.  Fisher regularity is checkable
from data: it requires only that the model's output distributions are
non-degenerate and smoothly parameterized.

\subsection{Open conjectures}

\Cref{tab:conjectures} summarizes the status of all conjectures in this branch.

\begin{table}[htbp]
\centering
\caption{Conjecture status summary}
\label{tab:conjectures}
\begin{tabular}{@{}llll@{}}
\toprule
\textbf{ID} & \textbf{Statement} & \textbf{Type} & \textbf{Status} \\
\midrule
\ref{conj:fisher-mixing} & Fisher gap governs mixing &
Central & Open (testable) \\
\ref{conj:attention-transport} & Attention as $\alpha$-transport &
Direction & Open (testable) \\
\ref{conj:pythagorean} & Pythagorean decomposition &
Direction & Open \\
\ref{conj:layer-norm} & Layer norm as Fisher normalization &
Direction & Open (testable) \\
\ref{conj:residual} & Residual stream as geodesic &
Direction & Open (testable) \\
\ref{conj:phase-transition} & Phase transitions as Fisher degeneracy &
Cross-branch & Finite-model testable \\
\ref{conj:sheaf-fisher} & Obstruction as Fisher incompatibility &
Cross-branch & Open \\
\ref{conj:data-processing} & Data processing and $C$-theorem &
Cross-branch & Open \\
\ref{conj:fisher-localization} & Fisher-weighted localization &
Cross-branch & Finite-model testable \\
\bottomrule
\end{tabular}
\end{table}

For a journal submission, Conjecture~\ref{conj:fisher-mixing} is the central
open problem.  Conjectures~\ref{conj:phase-transition}
and~\ref{conj:fisher-localization} are finite-model testable and could be
resolved computationally in subsequent work.  All others are directions for
future research.

\subsection{Cross-branch synthesis summary}

Table~\ref{tab:cross-branch} summarizes how each preceding branch acquires a
Fisher-metric interpretation under the identification of Branch~X.

\begin{table}[htbp]
\centering
\caption{\textbf{Cross-branch synthesis.}  Each row shows a parent-branch
result (left) and its Fisher-metric interpretation (right).  Entries marked
with $\star$ are conjectural.}
\label{tab:cross-branch}
\begin{tabularx}{\linewidth}{@{}p{0.12\linewidth}p{0.36\linewidth}p{0.44\linewidth}@{}}
\toprule
\textbf{Branch} & \textbf{Parent result} & \textbf{Fisher-metric interpretation} \\
\midrule
VI &
Free energy $\psi = \log Z$; susceptibility $\chi$ &
$\gF = \nabla^2\psi$; Fisher metric = susceptibility (Thm.~\ref{thm:fisher-free-energy}) \\
\addlinespace[2pt]
VI &
Completion transition at $\Gamma = 1$ &
Fisher degeneracy $\det(\gF) \to 0$ at transition$^\star$ (Conj.~\ref{conj:phase-transition}) \\
\addlinespace[2pt]
VII &
\v{C}ech obstruction $\delta(b) \in \check{H}^1$ &
Failure of local Fisher metrics to glue globally$^\star$ (Conj.~\ref{conj:sheaf-fisher}) \\
\addlinespace[2pt]
VIII &
$C$-theorem: entropy monotone under RG &
Data processing inequality for Fisher information$^\star$ (Conj.~\ref{conj:data-processing}) \\
\addlinespace[2pt]
IX &
Spectral gap $\mu_1$ controls capacity &
Fisher Laplacian gap $\mu_1^{\mathrm{F}}$ (Thm.~\ref{thm:fisher-csp}); sandwiching by $\kappa_{\mathrm{F}}$ \\
\addlinespace[2pt]
IX &
Completion count $N_k$ localizes in eigenspace &
Fisher-weighted $\tilde{N}_k^{\mathrm{F}}$ amplifies signal$^\star$ (Conj.~\ref{conj:fisher-localization}) \\
\bottomrule
\end{tabularx}
\end{table}

\subsection{Instantiation: Fisher-metric objects in running examples}

Table~\ref{tab:fisher-instantiation} shows how the abstract Fisher-metric
objects of Branch~X instantiate in each running example, paralleling the
instantiation tables in \cite{davis2025manifold} (Tables~3--4) and
\cite{davis2025sameness} (Section~7).

\begin{table}[htbp]
\centering
\caption{\textbf{Fisher-metric instantiation in running examples.}}
\label{tab:fisher-instantiation}
\begin{adjustbox}{max width=\linewidth}
\begin{tabular}{p{0.22\linewidth}p{0.24\linewidth}p{0.24\linewidth}p{0.24\linewidth}}
\toprule
\textbf{Object} & \textbf{HERALD} & \textbf{VIDAR} & \textbf{AMC} \\
\midrule
Statistical manifold $\mathcal{S}$ &
Lineage posterior space &
vMF identity space on $\mathbb{S}^{d-1}$ &
Reconstruction-error distribution space \\
\addlinespace[2pt]
Fisher metric $\gF$ &
Hessian of log-partition of lineage model &
$\kappa A_d(\kappa)(\delta_{ij} - \mu_i\mu_j)$ &
Sensitivity of format translation to state perturbation \\
\addlinespace[2pt]
Exponential family? &
Approximately (Boltzmann over lineages) &
Yes (von Mises--Fisher) &
Approximately (Gaussian reconstruction error) \\
\addlinespace[2pt]
Natural params $\theta$ &
Log-odds of lineage assignment &
$\kappa\mu$ (concentration $\times$ direction) &
Mean reconstruction error per format \\
\addlinespace[2pt]
Expectation params $\eta$ &
Lineage posterior means &
$A_d(\kappa)\mu$ (mean direction) &
Expected reconstruction quality \\
\addlinespace[2pt]
Dual connections &
$\nabla^{(e)}$: flat in log-odds; $\nabla^{(m)}$: flat in posteriors &
$\nabla^{(e)}$: flat in $\kappa\mu$; $\nabla^{(m)}$: flat in mean direction &
$\nabla^{(e)}$: flat in log-error; $\nabla^{(m)}$: flat in expected error \\
\addlinespace[2pt]
Measurable via &
Assay posterior logits &
Softmax logits at output layer &
Translation loss gradients \\
\bottomrule
\end{tabular}
\end{adjustbox}
\end{table}

\subsection{Result count and framework tracking}

Table~\ref{tab:result-count} summarizes the formal results in this branch.

\begin{table}[htbp]
\centering
\caption{\textbf{Branch~X result count.}  The ``Journal paper'' column
indicates which results would appear in a standalone journal submission
(see Strategic Note in the outline document).}
\label{tab:result-count}
\begin{tabular}{lcc}
\toprule
\textbf{Result type} & \textbf{Branch document} & \textbf{Journal paper} \\
\midrule
Definitions & 5 & 3 \\
Theorems & 10 & 5 \\
Lemmas & 2 & 1 \\
Corollaries & 5 & 2 \\
Propositions & 2 & 1 \\
Remarks & 5 & 1 \\
Conjectures & 9 & 1 (central) \\
Algorithms & 2 & 2 \\
\midrule
\textbf{Total formal results} & \textbf{40} & $\sim$\textbf{15} \\
\bottomrule
\end{tabular}
\end{table}

\paragraph{Running framework total.}
With Branch~X, the Davis framework contains approximately $252 + 40 = 292$
formal results across Branches~I--X.  The information-geometric identification
provides the first principled metric-selection mechanism and the first
empirically testable measurement program for the framework's curvature bounds.


%====================================================
% APPENDIX
%====================================================
\appendix
\section{Detailed proofs}\label{sec:appendix}

\subsection{Supplementary details for Theorem~\ref{thm:curvature-capacity}
(Curvature capacity from Cram\'er--Rao)}
\label{subsec:proof-curvature-capacity}

The complete proof---including the $d$-dimensional Bishop--Gromov volume
packing argument that contains the parent framework's $C \lesssim 1/K$
law as the $d{=}2$ specialization---is given in Section~\ref{sec:curvature-capacity}.  Here we record the
distinguishability-theoretic foundation that connects the $\delta$-packing
definition to statistical hypothesis testing.

\paragraph{Distinguishability via Pinsker.}
Two distributions $p_{\gamma(s)}$ and
$p_{\gamma(t)}$ along a completion path $\gamma$ are distinguishable at error level $\delta$
if and only if there exists a test $\phi : \mathcal{X} \to \{0,1\}$ with
\[
\Prob_{p_{\gamma(s)}}(\phi = 1) + \Prob_{p_{\gamma(t)}}(\phi = 0) \le \delta.
\]
By the Neyman--Pearson lemma, optimal distinguishability is controlled by the
total variation distance $\|p_{\gamma(s)} - p_{\gamma(t)}\|_{\mathrm{TV}}$,
which is bounded below by the Fisher-Rao distance via Pinsker's inequality:
\begin{equation}\label{eq:pinsker}
\|p_s - p_t\|_{\mathrm{TV}} \ge \frac{1}{\sqrt{2}}\,
d_{\mathrm{F}}(\gamma(s), \gamma(t)),
\end{equation}
where $d_{\mathrm{F}}$ is the geodesic distance in $(\mathcal{S}, \gF)$.
This justifies the identification of ``distinguishable completions'' with
``$\delta$-separated points in the Fisher-Rao metric'' that underlies the
packing argument in Theorem~\ref{thm:curvature-capacity}.

\subsection{Proof of Theorem~\ref{thm:fisher-spectral-sandwich}
(Fisher spectral capacity sandwiching)}
\label{subsec:proof-spectral-sandwich}

\begin{proof}
Each bound is a specialization of a classical spectral-geometric inequality
to the specific metric $\gF$.

\paragraph{(i) Cheeger bound.}
The Cheeger constant of $(\mathcal{S}, \gF)$ is
$h_{\mathrm{F}} := \inf_{\Omega} \mathrm{Area}_{\gF}(\partial\Omega) /
\min\{\mathrm{Vol}_{\gF}(\Omega),
\mathrm{Vol}_{\gF}(\mathcal{S} \setminus \Omega)\}$,
where the infimum is over all smooth domains $\Omega \subset \mathcal{S}$.
Cheeger's inequality~\cite{cheeger1970} gives
$\mu_1 \ge h^2/4$ for any compact Riemannian manifold; the factor $1/2$
in the stated bound uses the improved constant from
Buser~\cite{buser1982}.  Specializing to $(\mathcal{S}, \gF)$ yields~(i).

\paragraph{(ii) Lichnerowicz bound.}
If $\mathrm{Ric}_{\gF} \ge (n-1)\kappa$ with $\kappa > 0$, then
Lichnerowicz's theorem~\cite{lichnerowicz1958} gives
$\mu_1 \ge n\kappa$.  Applied to $\gF$, this requires the Ricci curvature
of the Fisher metric to be positively bounded below---a non-trivial
condition that depends on the statistical manifold's geometry.  On compact
exponential families with bounded parameters, this is typically satisfied
with $\kappa$ depending on the parameter domain.

\paragraph{(iii) Cheng bound.}
Cheng's comparison theorem~\cite{cheng1975} gives
$\mu_1 \le (n-1)^2 \kappa / 4 + C/D^2$ on a manifold with Ricci curvature
$\ge (n-1)\kappa$ and diameter $D$.  When $\kappa \le 0$, this simplifies to
$\mu_1 \le C(n,\kappa)/D^2$.  Applied to $(\mathcal{S}, \gF)$ with
$D = D_{\mathrm{F}}$, this gives~(iii).

The Fisher spectral capacity
$C_{\mathrm{sp}}^{\mathrm{F}} = \mu_1^{\mathrm{F}} D_{\mathrm{F}}^2$
is thus bounded by the same Cheeger--Ricci--Cheng sandwich as the parent
framework's spectral capacity in Branch~IX, with all geometric quantities
computed in the Fisher metric.
\end{proof}

\subsection{Proof of Theorem~\ref{thm:fisher-csp}
(Fisher Laplacian on CSP graphs)}
\label{subsec:proof-fisher-csp}

\begin{proof}
We provide the full argument, including the Rayleigh quotient comparison.

The Fisher Laplacian is defined as
$\Delta^{\mathrm{F}}_G = D_{\mathrm{F}}^{-1/2}\, L_G\, D_{\mathrm{F}}^{-1/2}$,
where $D_{\mathrm{F}} = \diag(\gF(L_1), \dots, \gF(L_n))$ with
$\gF(L_k) = 1/N_k(L_k)$.  This is a symmetric matrix (since $L_G$ is
symmetric and $D_{\mathrm{F}}^{-1/2}$ is diagonal), and positive semidefinite
(since $L_G \succeq 0$ implies
$f^\top \Delta^{\mathrm{F}}_G f = (D_{\mathrm{F}}^{-1/2} f)^\top L_G
(D_{\mathrm{F}}^{-1/2} f) \ge 0$ for all $f$).

The null space of $\Delta^{\mathrm{F}}_G$ is
$\{f : D_{\mathrm{F}}^{-1/2} f \in \ker L_G\}
= \{D_{\mathrm{F}}^{1/2} \mathbf{1}\}$,
which is one-dimensional on connected graphs.  The first positive eigenvalue is
\[
\mu_1^{\mathrm{F}} = \min_{\substack{f \neq 0 \\ f \perp D_{\mathrm{F}}^{1/2}
\mathbf{1}}} \frac{f^\top \Delta^{\mathrm{F}}_G f}{f^\top f}.
\]
Substituting $h = D_{\mathrm{F}}^{-1/2} f$ and noting that
$f^\top f = h^\top D_{\mathrm{F}} h$, we get
\[
\mu_1^{\mathrm{F}} = \min_{\substack{h \neq 0 \\ h \perp \mathbf{1}}}
\frac{h^\top L_G h}{h^\top D_{\mathrm{F}} h}.
\]
Since $\gF_{\min} \|h\|^2 \le h^\top D_{\mathrm{F}} h \le \gF_{\max}
\|h\|^2$, we have
\[
\frac{h^\top L_G h}{\gF_{\max} \|h\|^2}
\le \frac{h^\top L_G h}{h^\top D_{\mathrm{F}} h}
\le \frac{h^\top L_G h}{\gF_{\min} \|h\|^2},
\]
and taking the minimum over $h \perp \mathbf{1}$ gives
\[
\frac{\mu_1(G)}{\gF_{\max}}
\le \mu_1^{\mathrm{F}}(G)
\le \frac{\mu_1(G)}{\gF_{\min}},
\]
which is~\eqref{eq:gap-relation}.  Here $\gF_{\max} = \max_k \gF(L_k)
= 1/\min_k N_k$ and $\gF_{\min} = \min_k \gF(L_k) = 1/\max_k N_k$.
\end{proof}

\subsection{Relationship between Amari--Chentsov tensor and parent curvature}
\label{subsec:ac-parent}

\begin{proposition}[Levi-Civita curvature from $\alpha$-curvatures]
\label{prop:lc-from-alpha}
The Riemann curvature tensor of the Levi-Civita connection $\nabla^{(0)}$ on
$(\mathcal{S}, \gF)$ is related to the $\alpha$-curvature tensors by
\begin{equation}\label{eq:lc-alpha-curvature}
R^{(0)}_{ijkl} = \frac{1}{2}\bigl(R^{(\alpha)}_{ijkl}
+ R^{(-\alpha)}_{ijkl}\bigr)
+ \frac{\alpha^2}{4}\bigl(T_{im} T_{jn} - T_{jm} T_{in}\bigr)
(\gF)^{mn} \gF_{kl},
\end{equation}
where $T_{ijk}$ is the Amari--Chentsov tensor
(Definition~\ref{def:ac-tensor}).  In particular, when $R^{(e)} = 0$
(exponential families), the Levi-Civita curvature is entirely determined by
the Amari--Chentsov tensor:
\[
R^{(0)}_{ijkl} = \frac{1}{4}\bigl(T_{im} T_{jn} - T_{jm} T_{in}\bigr)
(\gF)^{mn} \gF_{kl}.
\]
\end{proposition}

\begin{proof}
This follows from Lemma~\ref{lem:alpha-from-ac} and the standard relation
between Christoffel symbols and curvature.  The $\alpha$-connection's
Christoffel symbols differ from the Levi-Civita's by $-(\alpha/2)T_{ijk}$.
Computing the Riemann tensor $R^{(\alpha)}$ from
$\Gamma^{(\alpha)} = \Gamma^{(0)} - (\alpha/2)T$ and expanding to second
order in $\alpha$ gives the stated formula.  The details are in
Amari~\cite{amari1985}, Chapter~3.
\end{proof}

\noindent\emph{Significance for the Davis framework.}  On exponential families
(such as softmax outputs), the parent framework's curvature is entirely
``Amari--Chentsov curvature'': it arises from the third-order statistics of
the score function, not from any intrinsic geometric curvature of the
parameter space.  The parent's curvature capacity law
$C \le c/\sup K^{(0)}$ can therefore be expressed as a bound involving only
the skewness tensor $T$---the ``cubic part'' of the statistical geometry.
This connects the parent's capacity bounds to higher-order cumulants of the
output distribution, opening a path to estimation from finite samples.

%====================================================
% BIBLIOGRAPHY
%====================================================
\begin{thebibliography}{99}

\bibitem{amari1985}
Amari, S. (1985).
\emph{Differential-Geometrical Methods in Statistics.}
Lecture Notes in Statistics~28, Springer.

\bibitem{amari1998}
Amari, S. (1998).
Natural gradient works efficiently in learning.
\emph{Neural Computation} 10(2), 251--276.

\bibitem{amari2000}
Amari, S. \& Nagaoka, H. (2000).
\emph{Methods of Information Geometry.}
Translations of Mathematical Monographs~191, AMS/Oxford.

\bibitem{ay2017}
Ay, N., Jost, J., L\^e, H.V., \& Schwachh\"ofer, L. (2017).
\emph{Information Geometry.}
Ergebnisse der Mathematik~64, Springer.

\bibitem{campbell1986}
Campbell, L.L. (1986).
An extended \v{C}encov characterization of the information metric.
\emph{Proceedings of the AMS} 98(1), 135--141.

\bibitem{cencov1982}
\v{C}encov, N.N. (1982).
\emph{Statistical Decision Rules and Optimal Inference.}
Translations of Mathematical Monographs~53, AMS.

\bibitem{efron1975}
Efron, B. (1975).
Defining the curvature of a statistical problem (with applications to
second-order efficiency).
\emph{Annals of Statistics} 3(6), 1189--1217.

\bibitem{kunstner2019}
Kunstner, F., Hennig, P., \& Balles, L. (2019).
Limitations of the empirical Fisher approximation for natural gradient
descent.
\emph{NeurIPS 2019.}

\bibitem{martens2020}
Martens, J. (2020).
New insights and perspectives on the natural gradient method.
\emph{JMLR} 21(146), 1--76.

\bibitem{davis2025manifold}
Davis, B.R. (2025).
\emph{The Davis Manifold: Geometry-First Detection with Compositional
Error Budgets.}
Zenodo.
\url{https://doi.org/10.5281/zenodo.17642038}

\bibitem{davis2025sameness}
Davis, B.R. (2025).
\emph{The Geometry of Sameness: An $\varepsilon$-Equivalence of Translation
and Distance.}
Zenodo.

\bibitem{davis2025statmech}
Davis, B.R. (2025).
\emph{Statistical Mechanics of Semantic Fields} (Branch~VI).
Zenodo.

\bibitem{davis2025sheaf}
Davis, B.R. (2025).
\emph{Sheaf-Theoretic Completion Obstruction} (Branch~VII).
Zenodo.

\bibitem{davis2025rg}
Davis, B.R. (2025).
\emph{Renormalization Group Structure of Completion Fields} (Branch~VIII).
Zenodo.

\bibitem{davis2025spectral}
Davis, B.R. (2025).
\emph{Spectral Geometry of Constraint-Satisfaction Graphs} (Branch~IX).
Zenodo.

\bibitem{cheeger1970}
Cheeger, J. (1970).
A lower bound for the smallest eigenvalue of the Laplacian.
In \emph{Problems in Analysis}, Princeton University Press, 195--199.

\bibitem{buser1982}
Buser, P. (1982).
A note on the isoperimetric constant.
\emph{Annales scientifiques de l'\'Ecole Normale Sup\'erieure} 15(2), 213--230.

\bibitem{lichnerowicz1958}
Lichnerowicz, A. (1958).
\emph{G\'eom\'etrie des groupes de transformations.}
Dunod, Paris.

\bibitem{cheng1975}
Cheng, S.-Y. (1975).
Eigenvalue comparison theorems and its geometric applications.
\emph{Mathematische Zeitschrift} 143(3), 289--297.

\bibitem{rao1945}
Rao, C.R. (1945).
Information and the accuracy attainable in the estimation of statistical
parameters.
\emph{Bulletin of the Calcutta Mathematical Society} 37, 81--91.

\bibitem{bonnet1855}
Bonnet, O. (1855).
Sur quelques propri\'et\'es des lignes g\'eod\'esiques.
\emph{Comptes rendus de l'Acad\'emie des Sciences} 40, 1311--1313.

\bibitem{myers1941}
Myers, S.B. (1941).
Riemannian manifolds with positive mean curvature.
\emph{Duke Mathematical Journal} 8(2), 401--404.

\bibitem{diaconis1991}
Diaconis, P. \& Stroock, D. (1991).
Geometric bounds for eigenvalues of Markov chains.
\emph{Annals of Applied Probability} 1(1), 36--61.

\end{thebibliography}

\end{document}
